% !TEX root = ../main.tex
\chapter{Black--Scholes Formula and Risk-Neutral Valuation}\label{ch:bs}
\begin{paracol}{2}

\begin{sources}
\begin{tabularx}{\linewidth}{@{}l l L@{}}
\toprule
\textbf{Lecture} & \textbf{Speaker} & \textbf{Content}\\
\midrule
2024 L21 & V.~Strela & Update of the yield curve after the September 2024 Fed meeting; two-horse race (the bookie who sets odds from the pool); forwards, calls, puts; one-step replication of a forward and of a call; risk-neutral probability; lognormal dynamics, why $\dd W\sim\sqrt{\dd t}$, heuristic It\^o formula; replication argument and the Black--Scholes PDE; boundary conditions and the closed-form formula; the risk-neutral density; put--call parity on Bloomberg screens (AAPL, IBM) and the implied-volatility smile; the digital option and the call spread.\\
2013 L19 & V.~Strela & The same material (nearly the same slides), with the discrete one-step market in more detail (bond units $b$, the interest-rate quiz on the forward), suggestions for final-paper topics (heat equation, formula by expectation), an IBM option screen of March 2006, put--call parity, and the digital option as an interview question.\\
\bottomrule
\end{tabularx}

\smallskip
\textbf{Files.} Slides \file{mit18_642_f24_lec21.pdf} (26 slides) and 2013 slides \file{MIT18_S096F13_lecnote19.pdf}
(25 slides); no problem set, case study or code accompanies this lecture in either year.
\textbf{Not recorded.} 2024 L22, \emph{The Spectrum of Systematic Trading Strategies in Liquid Instruments}
(R.~Garon, Millennium Management), has no video and no slides: the guest did not grant permission to publish them, so there
is nothing to summarise here; the next recorded lecture is L23 (machine learning, \cref{ch:ml}).
\textbf{Exercises.} None of the course problem sets deals with this lecture; the exercises (each after the section it practises)
are all marked ``(not from the course)''.
\end{sources}

\section*{Motivation}
Chapter~\ref{ch:linalg} priced a call in a one-period market with two states and found that the real-world
probability of ``up'' did not matter (\cref{ex:call}). This chapter takes that idea to its limit. The
question is: \emph{what is the price today of a payoff $f(S_T)$ that is fixed by the price of a stock at a
future date $T$?} The answer, due to Black, Scholes and Merton (1973), is
\begin{itemize}
  \item \emph{model-dependent but preference-free}: once one specifies how the stock moves, the price is
        determined, and neither the expected return of the stock nor the risk appetite of investors enters;
  \item \emph{constructive}: the price is the cost of a self-financing portfolio in stock and cash that
        \emph{replicates} the payoff, and that portfolio is also the hedge;
  \item \emph{an expectation}: the price equals a discounted expectation of the payoff under a fictitious
        \emph{risk-neutral} measure $\Q$, under which every traded asset grows at the risk-free rate;
  \item \emph{a PDE}: the same price solves the Black--Scholes equation, a heat equation in disguise.
\end{itemize}
The lecture is deliberately placed \emph{before} the stochastic-calculus lectures
(\cref{ch:stochcalc,ch:sde}): It\^o's formula is introduced heuristically, in the way a physicist would expand
a function of a Brownian path, and the rigorous statements are deferred. Read this chapter as the
\emph{application} that motivates the formalism of the next ones, and as the continuous-time counterpart of
\cref{sec:oneperiod}. Sections~\ref{ch15:sec-bookie}--\ref{ch15:sec-onestep} are the discrete world,
\cref{ch15:sec-tree,ch15:sec-lognormal} the passage to continuous time,
\cref{ch15:sec-pde,ch15:sec-rn,ch15:sec-formula} the two routes to the formula, and
\cref{ch15:sec-parity,ch15:sec-digital,ch15:sec-greeks,ch15:sec-smile} what one does with it in practice.

% ==================================================================
\section{Where the lecture starts: the yield curve}\label{ch15:sec-curve}
As promised in the first lecture, the class opens with an update of the US Treasury curve of \cref{fig:yieldcurve}.
Slide~2 overlays the curve of 3 September 2024 (just before the course began) on that of 15 November 2024
\ts{24}{21}{0:00}. The Fed had cut in the meantime by 75 basis points in total (in the calendar: 50~bp in September and
25~bp in November, a detail not spoken in class), and the short end of the curve moved down by about that
amount while the long end moved by only about 50~bp \ts{24}{21}{1:07}: not a parallel shift but a steepening.

The lecturer then puts the current shape next to three historical snapshots \ts{24}{21}{2:14}: the
curves of 2007 (just before the crisis) and 2019 were both \emph{inverted} and a recession followed, while by
September 2021 the curve had a normal shape and the recession had been short. He calls this ``non-scientific
prediction'': a serious analysis would decompose the curves with principal components (\cref{ch:pca}), and the
term structure is, in his words, in the centre of finance because the discount factors of \cref{sec:discounting}
are read off it. For the present chapter only one consequence matters: discounting uses a rate $r$, which the
lecture takes \emph{constant}. Everything that follows is an exercise in a single-rate world; making $r$ stochastic
is the subject of the short-rate and HJM models of \cref{ch:hjm}.

% ==================================================================
\section{Setting the odds: a bookie's version of hedging}\label{ch15:sec-bookie}
The lecture introduces risk-neutral valuation with an example that has no mathematics in it
\ts{24}{21}{4:21}\ts{13}{19}{0:00}. A bookie knows that horse~1 wins a race with probability $20\%$ and horse~2 with
probability $80\%$. The public does not know this, and bets \$10\,000 on horse~1 and \$50\,000 on horse~2
\ts{24}{21}{5:24}\ts{13}{19}{1:01}; the pool is \$60\,000 (a commission, say $1\%$, is what the bookie keeps in
the end and is ignored below). The bookie must announce the odds ``$k$ to one'': a winning stake of $s$ returns
$s+ks$.

\paragraph{Odds from the ``true'' probabilities: $4$ to $1$.} Odds $4:1$ equal the true ratio $0.8:0.2$. If horse~1
wins, the bookie collects \$60\,000 and pays $10\,000+4\cdot10\,000=\$50\,000$: profit \$10\,000. If horse~2 wins, he pays
$50\,000+\tfrac14\cdot50\,000=\$62\,500$: loss \$2\,500 \ts{24}{21}{6:28}. The expected profit is
$0.2\cdot10\,000+0.8\cdot(-2\,500)=0$, but only on average: the bookie carries the risk of the one race that is run.

\paragraph{Odds from the pool: $5$ to $1$.} Now ignore the true probabilities and set the odds by the money
staked, $50\,000/10\,000=5$. If horse~1 wins the bookie pays $10\,000+5\cdot10\,000=\$60\,000$; if horse~2 wins he pays
$50\,000+\tfrac15\cdot50\,000=\$60\,000$. Either way he pays exactly what he collected \ts{24}{21}{8:37}\ts{13}{19}{3:06}: the book is
\emph{hedged} and any commission is riskless profit. The pool proportions
\[
  q_1=\frac{10\,000}{60\,000}=\frac16,\qquad q_2=\frac{50\,000}{60\,000}=\frac56
\]
play the role of a \emph{risk-neutral} probability: whichever the true probabilities $(p_1,p_2)$, the bookie's
expected profit with odds set from $q$ is zero \emph{and its variance is zero}. The price of a bet on horse~1 (its stake,
relative to the payout) is $q_1=1/6$, not $p_1=1/5$.

\begin{intuition}[What the example teaches]
The bookie prices by \emph{what is already priced by the market} (the pool) and not by what he believes. The
numbers $q_i$ are not forecasts: they are the prices of the two Arrow--Debreu securities of \cref{def:onepmodel}
(``\$1 if horse $i$ wins'') normalised to sum to one. Three lessons carry over to derivatives \ts{24}{21}{9:42}:
\begin{itemize}
  \item an arbitrage-free price is the cost of a portfolio that reproduces the payoff in every outcome;
  \item that cost can be written as an expectation, but with the ``pool'' probabilities $q$ and not the
      true ones;
  \item the pool probabilities need not coincide with reality, and so a model of real-world
      probabilities is not needed to price and hedge.
\end{itemize}
\end{intuition}

\begin{coursenote}[Transcript slips]
In the 2013 transcript the bookie's payout on the second horse is given as ``\$12.25'' for the quarter of \$50\,000;
it is \$12\,500 (so the total payout is \$62\,500 and the loss \$2\,500, as on the slide). Also in 2013,
to the quiz ``which interest rate does the forward strike imply?'' a student answers ``2.5'', which the lecturer first
repeats and then corrects on the fly: $80\to88.41$ is about $10\%$ over the two years of \cref{ch15:sec-payoffs}, i.e.\ $5\%$ per year
continuously compounded, $80\,e^{0.05\cdot2}=88.41$ \ts{13}{19}{16:42}.
\end{coursenote}

\begin{exercise}[not from the course: the bookie]\label{ch15:ex-bookie}
Pools of $A$ on horse~1 and $B$ on horse~2. Show that odds set to $B/A$ to $1$ on horse~1 (and $A/B$ to $1$ on horse~2) give the bookie the same profit whichever horse wins, and identify the ``risk-neutral probabilities'' $q_i$.
Compute the standard deviation of the bookie's profit in the example if the odds are set from the true probabilities $(0.2,0.8)$, and compare with the pool odds.
\end{exercise}

% ==================================================================
\section{Forwards, calls and puts}\label{ch15:sec-payoffs}
Throughout, $S_t$ is the price at time $t$ of a stock (no dividends), $T$ the \emph{expiry} of the contract, $K$ its
\emph{strike} and $\tau=T-t$ the time to expiry. The three simplest derivatives pay, at $T$ \ts{24}{21}{10:47}\ts{13}{19}{4:06},
\[
  \text{forward: } S_T-K,\qquad \text{call: } (S_T-K)^+,\qquad \text{put: } (K-S_T)^+ ,
\]
with $x^+=\max(x,0)$. A \emph{forward contract} is an obligation to buy at $T$ at the price $K$ fixed today; $K$ is chosen so that
the contract costs nothing to enter \ts{24}{21}{11:54}\ts{13}{19}{5:07}. A \emph{call option} is the right, not the obligation, to
buy at $K$, and a \emph{put option} the right to sell at $K$ \ts{24}{21}{12:55}. The call is insurance against
the stock going up if one is short, and the put is insurance against a fall: its owner is protected below $K$. Unlike a forward,
an option costs a premium today, because it can only win. A call (put) is \emph{in the money} if $S>K$ ($S<K$), \emph{at the money} if $S=K$, and
\emph{out of the money} otherwise; the lecture always means \emph{European} options, exercisable only at $T$
(American options, exercisable any time, have no closed form and are computed on trees or finite-difference grids).

\begin{figure}[ht]
\centering
\begin{tikzpicture}
\begin{axis}[name=A,width=0.48\linewidth,height=6.2cm,xmin=0,xmax=200,ymin=-90,ymax=125,axis lines=left,
  xlabel={$S$},tick label style={font=\footnotesize},label style={font=\small},title={\small call and forward},
  title style={yshift=-3pt},legend style={at={(0.97,0.03)},anchor=south east,font=\scriptsize,draw=none,fill=none}]
  \addplot[thick,black!45] coordinates {(0,0) (80,0) (200,120)};
  \addplot[very thick,defcol] coordinates {(2,0.0000) (8,0.0000) (14,0.0000) (20,0.0000) (26,0.0004) (32,0.0077) (38,0.0574) (44,0.2499) (50,0.7561) (56,1.7730) (62,3.4588) (68,5.8914) (74,9.0629) (80,12.9014) (86,17.3009) (92,22.1466) (98,27.3322) (104,32.7679) (110,38.3825) (116,44.1224) (122,49.9485) (128,55.8331) (134,61.7571) (140,67.7071) (146,73.6744) (152,79.6530) (158,85.6391) (164,91.6300) (170,97.6241) (176,103.6203) (182,109.6178) (188,115.6161) (194,121.6151) (200,127.6144)};
  \addplot[very thick,dashed,thmcol] coordinates {(0,-88.41) (200,111.59)};
  \addplot[very thick,dotted,intcol] coordinates {(2,-77.9967) (8,-71.9967) (14,-65.9967) (20,-59.9967) (26,-53.9967) (32,-47.9967) (38,-41.9967) (44,-35.9967) (50,-29.9967) (56,-23.9967) (62,-17.9967) (68,-11.9967) (74,-5.9967) (80,0.0033) (86,6.0033) (92,12.0033) (98,18.0033) (104,24.0033) (110,30.0033) (116,36.0033) (122,42.0033) (128,48.0033) (134,54.0033) (140,60.0033) (146,66.0033) (152,72.0033) (158,78.0033) (164,84.0033) (170,90.0033) (176,96.0033) (182,102.0033) (188,108.0033) (194,114.0033) (200,120.0033)};
  \legend{payoff $(S-80)^+$,{call price $C(t,S)$},forward payoff $S-K$,{forward price $F(t,S)$}}
\end{axis}
\begin{axis}[at={(A.east)},anchor=west,xshift=1.1cm,width=0.48\linewidth,height=5.6cm,xmin=0,xmax=200,ymin=-5,ymax=80,axis lines=left,
  xlabel={$S$},tick label style={font=\footnotesize},label style={font=\small},title={\small put},title style={yshift=-3pt},
  legend style={at={(0.97,0.97)},anchor=north east,font=\scriptsize,draw=none,fill=none}]
  \addplot[thick,black!45] coordinates {(0,80) (80,0) (200,0)};
  \addplot[very thick,defcol] coordinates {(2,70.3870) (8,64.3870) (14,58.3870) (20,52.3870) (26,46.3874) (32,40.3947) (38,34.4444) (44,28.6369) (50,23.1431) (56,18.1600) (62,13.8458) (68,10.2784) (74,7.4499) (80,5.2884) (86,3.6879) (92,2.5336) (98,1.7192) (104,1.1549) (110,0.7694) (116,0.5094) (122,0.3355) (128,0.2201) (134,0.1441) (140,0.0941) (146,0.0614) (152,0.0400) (158,0.0261) (164,0.0170) (170,0.0111) (176,0.0073) (182,0.0048) (188,0.0031) (194,0.0021) (200,0.0014)};
  \legend{payoff $(80-S)^+$,{put price $P(t,S)$}}
\end{axis}
\end{tikzpicture}
\caption{Payoffs at expiry (grey or dashed) and prices two years before expiry (coloured), as functions of the stock price,
for $K=80$, $r=5\%$, $\sigma=20\%$, $\tau=2$ years (call, put) and $K=88.41=80e^{0.1}$ (forward), the numbers of slides~3--6.
The volatility is not stated on the slides; $20\%$ reproduces the digital curve of slide~26 (\cref{ch15:sec-digital}).
The price curves are the Black--Scholes formula derived below; the forward price is $S-Ke^{-r\tau}$.}\label{ch15:fig-payoffs}
\end{figure}

\Cref{ch15:fig-payoffs} anticipates the goal of the lecture \ts{24}{21}{13:55}: to compute the curved lines, the prices today
of the options, and the straight line for the forward. It also states the three messages of the lecture, which are exactly
the content of the box below \ts{24}{21}{15:01}\ts{13}{19}{9:22}.

\begin{keyformulas}
\textbf{Three messages of the lecture} (slide~7).
\begin{itemize}
  \item Given the current price of the stock \emph{and a model for its dynamics}, the price of a derivative is
        not uncertain.
  \item The price does not depend on the risk preferences of market participants (nor on the expected return of the stock).
  \item The mathematics of stochastic processes computes the price \emph{and the risks} (hedge ratios) of the derivative.
\end{itemize}
\end{keyformulas}

% ==================================================================
\section{One step: replication and the risk-neutral probability}\label{ch15:sec-onestep}
The lecture builds the theory in a market with one period of length $\dd t$ and two states \ts{24}{21}{16:04}\ts{13}{19}{11:27}. It is
the market of \cref{sec:binomial} with continuous compounding: the stock moves from $S_0$ to $S_1$ (\emph{up}, with real-world
probability $p$) or $S_2<S_1$ (\emph{down}), a riskless account grows from $B$ to $Be^{r\dd t}$, and a derivative $f$ pays $f_1$ or $f_2$.

\subsection{The forward}\label{ch15:sec-fwd1}
A forward on the stock pays $S-K_0$ at $\dd t$ and costs nothing today. \emph{The naive approach} sets the price
equal to the expected payoff with the real-world $p$,
$p(S_1-K_0)+(1-p)(S_2-K_0)=0$, i.e.\ $K_0=pS_1+(1-p)S_2$ (equal to $\frac12(S_1+S_2)$ if $p=\frac12$)
\ts{24}{21}{17:05}\ts{13}{19}{13:31}. It is wrong, because nobody knows $p$, and it is exploitable, as follows
\ts{24}{21}{19:08}\ts{13}{19}{14:33}.

\begin{proposition}[Forward price]\label{ch15:prop-fwd}
In the absence of arbitrage, $K_0=S_0e^{r\dd t}$: the strike of a zero-cost forward is the price of the stock capitalised at the
risk-free rate, whatever the real-world dynamics.
\end{proposition}
\begin{proof}
At time~$0$, borrow $S_0$ from the bank, buy the stock and sell the forward with strike $K_0$: the net cost is zero. At $\dd t$ deliver
the stock in exchange for $K_0$ and repay $S_0e^{r\dd t}$: the payoff is $K_0-S_0e^{r\dd t}$ in both states, a certain
amount that cost nothing. If it were positive this would be an arbitrage; if negative, the opposite position would
be \ts{24}{21}{20:11}\ts{13}{19}{15:35}. Hence $K_0=S_0e^{r\dd t}$.
\end{proof}
So the stock itself (financed with a loan) \emph{replicates} the forward: it has exactly the same payoff, and no assumption on $p$, $S_1$, $S_2$
or on how the stock ``really'' behaves was used. In the lecture's example, $S_0=80$, $T=2$ years and $K=88.41$: $88.41/80=e^{0.1}$, so $r=5\%$
continuously compounded \ts{13}{19}{16:42}. With $r=0$ one gets $K_0=S_0$, and one can still ask for a probability that gives the
same price as the expectation: the number $q$ with $K_0=qS_1+(1-q)S_2$ satisfies
\begin{equation}\label{ch15:eq-q}
  q=\frac{S_0e^{r\dd t}-S_2}{S_1-S_2}\in(0,1)\quad\text{iff}\quad S_2<S_0e^{r\dd t}<S_1,\qquad
  qS_1+(1-q)S_2=S_0e^{r\dd t}.
\end{equation}
This $q$ (called $p_n$ or $p_{rn}$ on the slides, $q$ in 2013) is the \emph{risk-neutral probability}: under it the stock earns exactly the
risk-free rate on average \ts{24}{21}{23:25}. The condition $0<q<1$ is the no-arbitrage condition of \cref{prop:binomial-na}
with $1+r_fT$ replaced by $e^{r\dd t}$.

\subsection{A general claim}\label{ch15:sec-claim1}
For a call (or any claim) the stock alone is no longer enough, but a mix of stock and cash is. Look for a stock position $a$ (in units of stock)
and a cash amount $B$ (in the riskless account, growing to $Be^{r\dd t}$) with \ts{24}{21}{25:30}--\ts{24}{21}{28:33}\ts{13}{19}{19:51}
\begin{equation}\label{ch15:eq-repl}
  f_1=aS_1+Be^{r\dd t},\qquad f_2=aS_2+Be^{r\dd t}.
\end{equation}
Subtracting, $a=\dfrac{f_1-f_2}{S_1-S_2}$, and then $B=e^{-r\dd t}\dfrac{S_1f_2-S_2f_1}{S_1-S_2}$. By no arbitrage $f_0=aS_0+B$
\ts{24}{21}{29:42}\ts{13}{19}{20:55}, and substituting one finds
\begin{equation}\label{ch15:eq-f0}
  f_0=e^{-r\dd t}\bigl[q\,f_1+(1-q)\,f_2\bigr],\qquad q=\frac{S_0e^{r\dd t}-S_2}{S_1-S_2},
\end{equation}
the same $q$ as in \eqref{ch15:eq-q}. Indeed
$aS_0+B=e^{-r\dd t}\bigl[\frac{(f_1-f_2)S_0e^{r\dd t}+S_1f_2-S_2f_1}{S_1-S_2}\bigr]
=e^{-r\dd t}\bigl[f_1\frac{S_0e^{r\dd t}-S_2}{S_1-S_2}+f_2\frac{S_1-S_0e^{r\dd t}}{S_1-S_2}\bigr]$
\ts{24}{21}{31:51}\ts{13}{19}{21:59}. This is \eqref{eq:rnbinomial} of \cref{sec:binomial} again.

\begin{proposition}[One-step risk-neutral pricing]\label{ch15:prop-onestep}
In the one-step market with $S_2<S_0e^{r\dd t}<S_1$, every claim $f$ is replicable, and its unique arbitrage-free price is
$f_0=e^{-r\dd t}\E_\Q[f_1]$ with $\Q=(q,1-q)$ given by \eqref{ch15:eq-q}. Equivalently, both $e^{-rt}S_t$ and $e^{-rt}f_t$ are martingales under $\Q$
(a two-point statement of the first fundamental theorem, \cref{thm:ftap}).
\end{proposition}
The price of the call is thus computed by taking \emph{expectations of its payoff}, $C_0=e^{-r\dd t}q(S_1-K)$ if $S_2<K<S_1$
(slide~12) \ts{24}{21}{31:51}. The bookie's pool proportions play the role of $q$; a physicist would say that $\Q$ is the measure that makes the
\emph{discounted} price a martingale, i.e.\ an unbiased random walk.

\begin{example}[Numbers of \cref{ex:call} in continuous compounding]\label{ch15:ex-call1}
Take $S_0=100$, $S_1=130$, $S_2=90$, $K=100$ and $e^{r\dd t}=1.03$ (the value used in \cref{ex:call}). Then $a=30/40=0.75$,
$B=-0.75\cdot90/1.03=-65.53$, $q=(103-90)/40=0.325$, and $f_0=0.75\cdot100-65.53=0.325\cdot30/1.03=9.466$, in agreement with
\cref{ex:call}. The real-world probability $p$ never appears.
\end{example}

\begin{intuition}[Replicating portfolio versus riskless portfolio]\label{ch15:int-riskless}
Equation~\eqref{ch15:eq-repl} says that the cash-plus-stock portfolio $aS+B$ has the same value as $f$ in both states. Turn it around
\ts{24}{21}{32:52}: the portfolio ``one unit of the derivative, minus $a$ units of stock'' pays $f_i-aS_i=Be^{r\dd t}$ in both
states: it is \emph{riskless}, and must therefore grow at the risk-free rate. For a physicist this is a statement that the hedged
combination has zero ``noise'', i.e.\ $a=\Delta f/\Delta S$ is a discrete derivative that cancels the fluctuation. In continuous time the same argument is
used in the form ``riskless portfolio'' (\cref{ch15:sec-pde}). Which of the two forms is more convenient depends on the problem \ts{24}{21}{32:52}.
\end{intuition}

\begin{exercise}[not from the course: one-step market]\label{ch15:ex-onestep}
$S_0=50$, $S_1=60$, $S_2=42$, $r\dd t=\ln1.02$.
\begin{enumerate}[(a)]
  \item Compute $q$, check the no-arbitrage condition and find the fair forward strike.
  \item Price a call ($K=50$) and a put ($K=50$) by replication and by $\E_\Q$ and verify parity.
  \item Let the real-world probability of ``up'' be $p=0.9$: compute $e^{-r\dd t}\E_{\Prob}[\text{call}]$ and the profit of a trader who sells the call at that price and hedges with $a$ shares.
\end{enumerate}
\end{exercise}


% ==================================================================
\section{Many steps: the binomial tree and its limit}\label{ch15:sec-tree}
This section is not spelled out in the lecture, but it is the bridge between the one-step market of
\cref{ch15:sec-onestep} and the continuous model of the slides (the second bullet of 2013 slide~21, ``tree methods, equivalent to
explicit scheme''); it is also the way most exotic options are still priced.

Split $[0,T]$ into $n$ steps of length $\dd t=T/n$. In each step the stock is multiplied by $u$ or by $d$; the
Cox--Ross--Rubinstein (CRR) choice is
\begin{equation}\label{ch15:eq-crr}
  u=e^{\sigma\sqrt{\dd t}},\qquad d=\frac1u=e^{-\sigma\sqrt{\dd t}},\qquad
  q=\frac{e^{r\dd t}-d}{u-d}.
\end{equation}
Each node of the tree is a copy of the one-step market, so \cref{ch15:prop-onestep} applies at every node and prices are computed by
\emph{backward induction} from the payoff at the leaves,
\begin{equation}\label{ch15:eq-backward}
  f(S,t)=e^{-r\dd t}\bigl[q\,f(uS,t+\dd t)+(1-q)\,f(dS,t+\dd t)\bigr],
  \qquad f(S,T)=\text{payoff}(S).
\end{equation}
Unrolling it, $f(S_0,0)=e^{-rT}\E_\Q[\text{payoff}(S_T)]$ with $S_T=S_0u^{J}d^{n-J}$ and $J\sim\mathrm{Bin}(n,q)$. Since
the number $q$ was fixed by the requirement that the \emph{discounted} stock be a martingale, this is the multi-period FTAP of \cref{thm:ftap}
and the martingale $e^{-rt}S_t$ of \cref{ch:sp1}.

\begin{proposition}[Two limits of the tree]\label{ch15:prop-limit}
Let $\dd t\to0$ with \eqref{ch15:eq-crr}.
\begin{enumerate}[(i)]
  \item \emph{Distribution.} $q=\frac12+\frac{r-\sigma^2/2}{2\sigma}\sqrt{\dd t}+O(\dd t)$ and
        $\ln(S_T/S_0)\Rightarrow N\bigl((r-\tfrac12\sigma^2)T,\ \sigma^2T\bigr)$. The price of a call
        converges to $e^{-rT}\E_\Q[(S_T-K)^+]$ with $S_T$ lognormal, i.e.\ to the Black--Scholes value of \cref{ch15:sec-formula}.
  \item \emph{Equation.} If $f$ is smooth, \eqref{ch15:eq-backward} implies
        $f_t+\tfrac12\sigma^2S^2f_{SS}+rSf_S-rf=o(1)$, the Black--Scholes PDE of \cref{ch15:sec-pde}.
\end{enumerate}
\end{proposition}
\begin{proof}
(i) Expand $u=1+\sigma\sqrt{\dd t}+\tfrac12\sigma^2\dd t+O(\dd t^{3/2})$, $d=1-\sigma\sqrt{\dd t}+\tfrac12\sigma^2\dd t+O(\dd t^{3/2})$, so
$u-d=2\sigma\sqrt{\dd t}\,(1+O(\dd t))$ and $e^{r\dd t}-d=\sigma\sqrt{\dd t}+(r-\tfrac12\sigma^2)\dd t+O(\dd t^{3/2})$, which gives the
expansion of $q$. A step of $\ln S$ is $\sigma\sqrt{\dd t}\,\xi$ with $\xi=\pm1$, $\Prob_\Q(\xi=1)=q$; its mean is
$\sigma\sqrt{\dd t}(2q-1)=(r-\tfrac12\sigma^2)\dd t+O(\dd t^{3/2})$ and its variance $\sigma^2\dd t\,[1-(2q-1)^2]=\sigma^2\dd t\,(1+O(\dd t))$.
Summing $n=T/\dd t$ independent steps, the mean tends to $(r-\tfrac12\sigma^2)T$ and the variance to $\sigma^2T$; the steps are bounded by $\sigma\sqrt{\dd t}\to0$,
so the Lindeberg central limit theorem (\cref{ch:probability}) applies. The payoff is continuous with linear growth and
$\E_\Q[S_T]=S_0e^{rT}$ exactly for every $n$, which supplies the uniform integrability needed to pass to the limit in the expectation.

(ii) By the definition of $q$, $q(u-1)+(1-q)(d-1)=e^{r\dd t}-1$ exactly, and $q(u-1)^2+(1-q)(d-1)^2=\sigma^2\dd t\,(1+O(\sqrt{\dd t}))$. Taylor-expanding
$f(uS,t+\dd t)$ and $f(dS,t+\dd t)$ to second order in $(u-1)S$, $(d-1)S$ and first order in $\dd t$ and inserting in \eqref{ch15:eq-backward} gives
$e^{r\dd t}f=f+f_t\dd t+(e^{r\dd t}-1)Sf_S+\tfrac12\sigma^2S^2f_{SS}\dd t+o(\dd t)$; subtract $f$ and divide by $\dd t$.
\end{proof}

\Cref{ch15:tab-crr} and \cref{ch15:fig-crr} show the convergence for the option of \cref{ch15:fig-payoffs} at the money:
$S_0=K=80$, $T=2$, $r=5\%$, $\sigma=20\%$, for which the formula of \cref{ch15:sec-formula} gives $C=12.9014$. The error decays like $1/n$ but alternates in sign
with the parity of $n$, because at the money the strike sits on a node for even $n$ and in the middle of a lattice cell for odd $n$: the column $n\cdot\text{error}$ is
stable near $-2.26$ for even $n$ and near $+1.73$ for odd $n$ (e.g.\ $n=51,101,1001$), so the oscillation does not die out, only its amplitude shrinks. The single-step value $14.51$ is more than $12\%$ off.

\begin{table}[ht]
\centering\small
\begin{tabular}{r r r r r}
\toprule
$n$ & tree price & error & $n\cdot$error & $q$\\
\midrule
1 & 14.5106 & $+1.6092$ & $+1.61$ & 0.6132\\
2 & 11.8731 & $-1.0283$ & $-2.06$ & 0.5775\\
5 & 13.2506 & $+0.3492$ & $+1.75$ & 0.5481\\
10 & 12.6786 & $-0.2228$ & $-2.23$ & 0.5338\\
50 & 12.8564 & $-0.0451$ & $-2.25$ & 0.5150\\
100 & 12.8789 & $-0.0226$ & $-2.26$ & 0.5106\\
1000 & 12.8992 & $-0.0023$ & $-2.26$ & 0.5034\\
5000 & 12.9010 & $-0.0005$ & $-2.27$ & 0.5015\\
\bottomrule
\end{tabular}
\caption{CRR tree price of the European call $S_0=K=80$, $T=2$, $r=5\%$, $\sigma=20\%$ (Black--Scholes: 12.9014). Computed for these notes.}\label{ch15:tab-crr}
\end{table}

\begin{figure}[ht]
\centering
\begin{tikzpicture}
\begin{axis}[width=0.85\linewidth,height=5.4cm,axis lines=left,xmin=0,xmax=60,ymin=11.5,ymax=14.7,
  xlabel={number of steps $n$},ylabel={price},tick label style={font=\small},label style={font=\small},
  legend style={at={(0.97,0.97)},anchor=north east,font=\scriptsize,draw=none,fill=none}]
  \addplot[only marks,mark=*,mark size=1.3pt,defcol] coordinates {(1,14.5106) (2,11.8731) (3,13.4811) (4,12.3590) (5,13.2506) (6,12.5342) (7,13.1508) (8,12.6240) (9,13.0952) (10,12.6786) (11,13.0599) (12,12.7153) (13,13.0354) (14,12.7416) (15,13.0175) (16,12.7614) (17,13.0038) (18,12.7768) (19,12.9930) (20,12.7892) (21,12.9842) (22,12.7993) (23,12.9770) (24,12.8078) (25,12.9710) (26,12.8150) (27,12.9658) (28,12.8211) (29,12.9614) (30,12.8264) (31,12.9575) (32,12.8311) (33,12.9541) (34,12.8352) (35,12.9511) (36,12.8389) (37,12.9484) (38,12.8422) (39,12.9460) (40,12.8451) (41,12.9438) (42,12.8478) (43,12.9418) (44,12.8502) (45,12.9400) (46,12.8524) (47,12.9384) (48,12.8545) (49,12.9369) (50,12.8564) (51,12.9355) (52,12.8581) (53,12.9342) (54,12.8597) (55,12.9330) (56,12.8612) (57,12.9319) (58,12.8626) (59,12.9308) (60,12.8638)};
  \addplot[thick,thmcol,dashed] coordinates {(0,12.9014) (60,12.9014)};
  \legend{CRR tree,Black--Scholes}
\end{axis}
\end{tikzpicture}
\caption{Convergence of the binomial price to the Black--Scholes value (same parameters as \cref{ch15:tab-crr}); the alternation between odd and even $n$ is visible.}\label{ch15:fig-crr}
\end{figure}

\begin{intuition}[Why $\sigma\sqrt{\dd t}$ and where $-\frac12\sigma^2$ comes from]
The steps of the tree are not independent of the time step: if the space step were $\propto\dd t$ the process would freeze,
and if it were larger than $\sqrt{\dd t}$ the variance would explode. The unique scaling that leaves a non-trivial limit is the diffusive one,
$\text{step}\sim\sqrt{\dd t}$, as in Donsker's theorem (\cref{ch11:thm-donsker}); this is what the lecturer proves
in the next section by counting variances. The drift correction $-\tfrac12\sigma^2$ in the exponent is the same lognormal correction as in
\cref{ch04:rem-logprice}: the tree has $\E_\Q[S_T]=S_0e^{rT}$ (the martingale property is exact), but $\ln S_T$ has mean
$(r-\tfrac12\sigma^2)T$ because $\E[e^{X}]=e^{\E[X]+\Var X/2}$ for a normal variable.
\end{intuition}

\begin{exercise}[not from the course: the tree]\label{ch15:ex-tree}
\leavevmode
\begin{enumerate}[(a)]
  \item Show that the CRR tree prices a forward $(S_T-K)$ exactly, for every $n$.
  \item Write a program that prices the call of \cref{ch15:tab-crr} and reproduces the $1/n$ error and the odd--even oscillation; then choose the
      strike lattice-aligned ($K=S_0u^{k}$) and observe the effect.
  \item Show that with $u=e^{\sigma\sqrt{\dd t}+\nu\dd t}$, $d=e^{-\sigma\sqrt{\dd t}+\nu\dd t}$ (drift $\nu$ in the tree) one still converges to Black--Scholes provided $q$ is defined by the martingale condition.
\end{enumerate}
\end{exercise}

% ==================================================================
\section{Lognormal dynamics and a heuristic It\^o formula}\label{ch15:sec-lognormal}

\subsection{The model}
The lecture models the stock by the \emph{lognormal} (geometric Brownian motion, \cref{ch11:eq-gbm}) dynamics
\ts{24}{21}{33:56}\ts{13}{19}{22:04}
\begin{equation}\label{ch15:eq-gbm}
  \dd S_t=\mu S_t\dd t+\sigma S_t\dd W_t,
\end{equation}
where $\mu$ is the (real-world) drift, $\sigma$ the \emph{volatility}, and $\dd W$ the increment of a Brownian motion: normal, with
$\E[\dd] W=0$ and $\operatorname{sd}(\dd W)=\sqrt{\dd t}$, independent for disjoint time intervals \ts{13}{19}{23:04}. The proportional change $\dd S/S$
has a deterministic part and a noise part. The lecture says clearly that the choice is an \emph{example}: the method needs only that the stock
has some dynamics, and other choices (stochastic volatility, jumps) lead to other equations of the same type
\ts{24}{21}{33:56}.

\subsection{Why the standard deviation is $\sqrt{\dd t}$}
In the 2024 lecture Strela stops to answer a question that ``always bothers'' him \ts{24}{21}{36:07}. Take a random walk $W_n=\sum_{i=1}^n\xi_i$ with steps $\xi_i=\pm1$ with
probability $\frac12$, uncorrelated: $\E[W_n]=0$ (a martingale) and $\Var W_n=\sum_i\E[\xi_i^2]=n$
\ts{24}{21}{37:10}. Now put the $n$ steps in a finite time $T$, $\dd t=T/n$, and let each step have size $\pm\delta$:
$\Var W=n\delta^2=T\delta^2/\dd t$. For the limit $\dd t\to0$ to be neither $0$ nor $\infty$ one needs $\delta^2/\dd t$ to stay finite, i.e.\ $\delta\propto\sqrt{\dd t}$
\ts{24}{21}{40:26}. That is the content of Donsker's theorem (\cref{ch11:thm-donsker}); the same
argument fixed the size of the steps of the tree \eqref{ch15:eq-crr}.

\subsection{Taylor expansion with a $\sqrt{\dd t}$ noise}
Let $f(S,t)$ be smooth. In a deterministic world the second-order Taylor formula is
$\dd f=f_t\dd t+f_S\dd S+\tfrac12f_{tt}\dd t^2+f_{tS}\dd t\dd S+\tfrac12f_{SS}\dd S^2+\dots$ \ts{24}{21}{43:35}. With $\dd S$ from \eqref{ch15:eq-gbm} the terms have very
different sizes \ts{24}{21}{44:45}: $\dd t\dd S\sim\dd t^{3/2}$ and $\dd t^2$ are negligible, but
\[
  (\dd S)^2=\sigma^2S^2(\dd W)^2+O(\dd t^{3/2}),\qquad(\dd W)^2\sim\dd t,
\]
is of the \emph{same} order as $\dd t$ and must be kept. One writes $(\dd W)^2=\dd t$ (\cref{ch11:eq-dB2}) and obtains, to order $\dd t$,
\ts{24}{21}{45:45}
\begin{equation}\label{ch15:eq-ito}
  \dd f=\Bigl(f_t+\mu Sf_S+\tfrac12\sigma^2S^2f_{SS}\Bigr)\dd t+\sigma Sf_S\dd W,
\end{equation}
\emph{It\^o's formula} for a function of $S$ and $t$ (slide~15). The lecturer calls it a ``quick and dirty derivation'' \ts{24}{21}{47:46}; the
rigorous statement, with the It\^o integral and the proof, is in \cref{ch:stochcalc}.

\begin{coursenote}[A student's objection]
A student asks how $(\dd W)^2=\dd t$ can hold when $\dd W$ is random \ts{24}{21}{55:16}. The answer given (by the lecturer and by the instructor, P.~Kempthorne) is that
the notation is sloppy but reflects a theorem: the \emph{quadratic variation} of Brownian motion over $[0,t]$ is exactly $t$,
$\sum_j(\Delta_jW)^2\to t$ almost surely as the mesh goes to zero (\cref{ch11:thm-qv}) \ts{24}{21}{56:21}. Each squared increment is random, of mean $\dd t$
and standard deviation $\sqrt2\,\dd t$, but their \emph{sum} over $n\sim1/\dd t$ intervals has mean $t$ and standard deviation
$\sqrt{2t\,\dd t}\to0$: a law of large numbers. Inside an integral $(\dd W)^2$ can therefore be replaced by $\dd t$.
\end{coursenote}

\begin{intuition}[Physicist's reading]
Equation~\eqref{ch15:eq-ito} is the Taylor expansion of a function of a Brownian path, with the ``ordinary'' second-order term promoted to first order
because the path is fractal ($\Delta W\sim\sqrt{\Delta t}$). Its physical avatar is the Langevin equation with multiplicative noise interpreted in the It\^o
(left-point) sense: $\tfrac12\sigma^2S^2f_{SS}$ is the ``noise-induced drift'' that a Stratonovich convention would absorb in the definition of the drift
(\cref{ch:stochcalc}). The same extra term makes $\ln S$ drift at $\mu-\tfrac12\sigma^2$ rather than $\mu$: taking $f=\ln S$ in \eqref{ch15:eq-ito}
gives $\dd\ln S=(\mu-\tfrac12\sigma^2)\dd t+\sigma\dd W$, and hence \cref{ch11:eq-gbm},
$S_T=S_t\exp\bigl[(\mu-\tfrac12\sigma^2)\tau+\sigma(W_T-W_t)\bigr]$.
\end{intuition}

% ==================================================================
\section{The Black--Scholes equation}\label{ch15:sec-pde}
We now do in continuous time what \cref{ch15:sec-claim1} did in one step \ts{24}{21}{47:46}\ts{13}{19}{23:04}. \textbf{Assumptions} (implicit in the lecture): the
stock follows \eqref{ch15:eq-gbm} with \emph{constant} $\mu$, $\sigma$; the interest rate $r$ is constant and the same for borrowing and lending; the stock
pays no dividend; trading is continuous, without transaction costs, with unrestricted short selling and fractional positions; the claim is European.

\subsection{Derivation by replication}
Let a portfolio hold $a_t=a(S_t,t)$ shares of the stock and a cash amount $B_t$ in the money-market account, which grows as
$\dd B$ (due to interest) $=rB\dd t$ \ts{24}{21}{49:52}. Its value is $V=aS+B$. The portfolio is \emph{self-financing}: no money is added or withdrawn,
so the change of value is only due to price moves and interest,
\begin{equation}\label{ch15:eq-selffin}
  \dd V=a\,\dd S+rB\dd t=(a\mu S+rB)\dd t+a\sigma S\dd W .
\end{equation}
(On the slides: $\dd f=a\,\dd S+\dd B$; the lecture writes $\dd B=rB\dd t$ and $B=f-aS$, the same thing.) We want the portfolio to \emph{replicate} the claim:
$V_t=f(S_t,t)$ at all times, with $f(S,T)$ the payoff. Equate \eqref{ch15:eq-ito} and \eqref{ch15:eq-selffin} \ts{24}{21}{51:59}\ts{13}{19}{25:12}.

\emph{Noise terms} (coefficients of $\dd W$): $\sigma Sf_S=a\sigma S$, so
\begin{equation}\label{ch15:eq-hedge}
  a=\frac{\partial f}{\partial S}\quad(\text{the \emph{delta}})\ts{24}{21}{57:27}.
\end{equation}
\emph{Drift terms} (coefficients of $\dd t$), with $B=f-aS=f-Sf_S$ from $V=f$:
$f_t+\mu Sf_S+\tfrac12\sigma^2S^2f_{SS}=\mu Sf_S+r(f-Sf_S)$. The real-world drift $\mu$ cancels and one is left with
\begin{equation}\label{ch15:eq-bs}
  \boxed{\ \frac{\partial f}{\partial t}+\frac12\sigma^2S^2\frac{\partial^2f}{\partial S^2}+rS\frac{\partial f}{\partial S}-rf=0\ }
\end{equation}
\ts{24}{21}{58:33}\ts{13}{19}{26:20}, the \emph{Black--Scholes equation} (slide~17), a backward parabolic PDE with final condition $f(S,T)=\text{payoff}(S)$. It is the result of Black and Scholes (1973); Scholes and Merton received the 1997 Nobel prize in economics (Black had died in 1995), and all three were associated with MIT \ts{24}{21}{1:01:47}.
If the claim traded at any price other than $V_t$, buying the cheaper and selling the dearer of $f$ and $V$ would be an arbitrage, which is why $f=V$.

The lecture derives the same equation in the ``riskless portfolio'' form \ts{24}{21}{49:52}. Let $\Pi=f-aS$ with $a$ held fixed over the interval
$[t,t+\dd t]$. By \eqref{ch15:eq-ito}, $\dd\Pi=\dd f-a\dd S=(f_t+\tfrac12\sigma^2S^2f_{SS}+(f_S-a)\mu S)\dd t+(f_S-a)\sigma S\dd W$.
Choosing $a=f_S$ kills \emph{both} $\mu$ and the noise: $\dd\Pi=(f_t+\tfrac12\sigma^2S^2f_{SS})\dd t$ is deterministic, hence must equal $r\Pi\dd t=r(f-Sf_S)\dd t$ by no arbitrage,
which is \eqref{ch15:eq-bs} again. The delta-hedged position (long the option, short $f_S$ shares) earns the risk-free rate.

\begin{keyformulas}
\textbf{What the equation says} (slide~18; \ts{24}{21}{1:00:39}).
\begin{itemize}[nosep]
  \item Every tradable claim on $S$, whatever its payoff, satisfies \eqref{ch15:eq-bs}; the payoff enters only through the final (and boundary) conditions.
  \item The real-world drift $\mu$ does not appear: only $\sigma$ and $r$ do. (Compare \cref{sec:binomial}: the probability of ``up'' does not appear.)
  \item The replicating portfolio is also a \emph{hedging strategy}: to hedge a short position in the claim, hold $a=\partial f/\partial S$ shares and put the rest, $f-Sf_S$, in the bank.
        For a traded, liquid stock this is what a derivatives desk does all day \ts{13}{19}{28:31}.
  \item With another model for $S$ the equation changes (different coefficients, possibly more state variables), but the derivation is identical.
\end{itemize}
\end{keyformulas}

\subsection{Final and boundary conditions}
As for any PDE one needs conditions \ts{24}{21}{1:03:55}\ts{13}{19}{30:40}. The \emph{final} condition is the payoff, $C(S,T)=(S-K)^+$ for a call and $P(S,T)=(K-S)^+$ for a put.
The \emph{boundary} conditions follow from the economic sense of the options \ts{24}{21}{1:04:56}\ts{13}{19}{31:49}:
\begin{align}
  \text{call:}\quad & C(0,t)=0,\qquad C(S,t)\sim S-Ke^{-r(T-t)}\ \ (S\to\infty),\label{ch15:eq-bccall}\\
  \text{put:}\quad  & P(0,t)=Ke^{-r(T-t)},\qquad P(S,t)\to0\ \ (S\to\infty).\label{ch15:eq-bcput}
\end{align}
The stock price $S=0$ is absorbing for geometric Brownian motion (a stock worth zero stays worthless), so the call is worthless and the put is a
certain payment of $K$ at $T$, worth $Ke^{-r(T-t)}$ today. For $S\to\infty$ the call is exercised almost surely and is worth the forward, $S-Ke^{-r(T-t)}$.

\begin{coursenote}[Erratum: the call at infinity]
Slide~19 of 2024 (slide~18 of 2013) states $C(\infty,t)\cong S$. To leading order this is right, but the correct asymptotics is
\eqref{ch15:eq-bccall}: the call behaves as the \emph{forward} $S-Ke^{-r(T-t)}$, as the 2013 lecturer says in words (``discounted $S$ minus $K$'') \ts{13}{19}{31:49}. The put condition
$P(0,t)=Ke^{-r(T-t)}$ is correct on both sets of slides.
\end{coursenote}

\subsection{Reduction to the heat equation}\label{ch15:sec-heat}
Slide~18 states that ``by a change of variables'' \eqref{ch15:eq-bs} becomes the heat equation $u_\tau=u_{xx}$; the 2024 lecturer calls finding
the change of variables a good exercise \ts{24}{21}{1:02:51}, and in 2013 he proposed it as a final-paper topic \ts{13}{19}{30:40}. Here it is
(the slide's normalisation $u_\tau=u_{xx}$ differs from ours by a rescaling $\tau\to\sigma^2\tau/2$).

Put $\tau=T-t$ (so $\partial_t=-\partial_\tau$) and $y=\ln S$ (so $Sf_S=f_y$ and $S^2f_{SS}=f_{yy}-f_y$). Then \eqref{ch15:eq-bs} reads
\[
  f_\tau=\tfrac12\sigma^2f_{yy}+\bigl(r-\tfrac12\sigma^2\bigr)f_y-rf .
\]
Remove the discounting with $f=e^{-r\tau}w$, which gives $w_\tau=\tfrac12\sigma^2w_{yy}+\bigl(r-\tfrac12\sigma^2\bigr)w_y$, and remove the drift with a
\emph{Galilean shift}, $x=y+(r-\tfrac12\sigma^2)\tau$, $w(y,\tau)=u(x,\tau)$: since $w_\tau=u_\tau+(r-\tfrac12\sigma^2)u_x$, $w_y=u_x$, $w_{yy}=u_{xx}$, one gets
\begin{equation}\label{ch15:eq-heat}
\begin{gathered}
  f(S,t)=e^{-r\tau}\,u(x,\tau),\qquad x=\ln S+\bigl(r-\tfrac12\sigma^2\bigr)\tau,\\
  \frac{\partial u}{\partial\tau}=\frac{\sigma^2}2\frac{\partial^2u}{\partial x^2},\qquad u(x,0)=\text{payoff}(e^x).
\end{gathered}
\end{equation}
(One checks it on exponentials: $f=e^{-r\tau}e^{\alpha x+\sigma^2\alpha^2\tau/2}$ satisfies \eqref{ch15:eq-bs} for every $\alpha$; verified symbolically for these notes.)
By \cref{ch11:prop-heat} the solution with initial datum $u_0$ is $u(x,\tau)=\int p(y,\tau\mid x)\,u_0(y)\dd y$, the Gaussian kernel \eqref{ch11:eq-kernel}
with variance $\sigma^2\tau$: that is, $u(x,\tau)=\E\bigl[u_0(x+\sigma\sqrt\tau Z)\bigr]$ with $Z\sim N(0,1)$. Written in the original variables,
\begin{equation}\label{ch15:eq-pde-solution}
  f(S,t)=e^{-r\tau}\,\E\Bigl[\text{payoff}\bigl(Se^{(r-\sigma^2/2)\tau+\sigma\sqrt\tau Z}\bigr)\Bigr],
\end{equation}
which is exactly the risk-neutral expectation of \cref{ch15:sec-rn}: the two routes to the price (PDE and expectation) give the same answer, as they must (this is the
Feynman--Kac formula, \cref{ch:sde}).

\begin{intuition}[Diffusion of a payoff backwards in time]
The pricing PDE is the heat equation run \emph{backwards}: the payoff at $T$ is the ``initial temperature profile'' and the price today is that profile after
diffusing for a time $\tau$ with diffusivity $\sigma^2/2$, in the frame of the log-price moving with velocity $r-\tfrac12\sigma^2$ and with a decay $e^{-r\tau}$ (discounting).
The hockey stick smoothed by a Gaussian is precisely the curved price line of \cref{ch15:fig-payoffs}; the smoothing width
$\sigma\sqrt\tau$ is the size of the ``typical move'' of the log-price until expiry, so an option far from the strike is the payoff, and an option near
the strike is worth more than its intrinsic value because the price may still cross $K$ (the ``time value'' or convexity value).
\end{intuition}

\begin{exercise}[not from the course: the heat equation]\label{ch15:ex-heat}
\leavevmode
\begin{enumerate}[(a)]
  \item Verify the change of variables of \cref{ch15:sec-heat} by direct substitution.
  \item Solve $u_\tau=\frac12\sigma^2u_{xx}$ for $u(x,0)=\ind\{x>\ln K\}$ and show that $f=e^{-r\tau}\Phi(d_2)$ (the digital).
  \item The slide's normalisation is $u_\tau=u_{xx}$: find the rescaling of $\tau$.
\end{enumerate}
\end{exercise}


% ==================================================================
\section{Risk-neutral valuation}\label{ch15:sec-rn}
The second route to the price is the one the lecturer wants to stress: ``mostly about risk-neutral pricing, and Black--Scholes is an illustration of the method''
\ts{24}{21}{0:00}. In discrete time the price was $f_0=e^{-r\dd t}\E_\Q[f_1]$ with the special probability $q$. The continuous-time statement is
\ts{24}{21}{1:06:01}\ts{13}{19}{33:55}
\begin{equation}\label{ch15:eq-rnval}
  f_t=e^{-r(T-t)}\,\E_\Q\bigl[f_T\mid S_t\bigr],\qquad
  \text{where under }\Q:\quad \dd S=rS\dd t+\sigma S\dd W^{\Q},
\end{equation}
so that $S_t=e^{-r(T-t)}\E_\Q[S_T\mid S_t]$: under $\Q$ the \emph{discounted} stock price $e^{-rt}S_t$ is a martingale (slides 13, 21).

\paragraph{How to find $\Q$.} The lecturer's argument \ts{24}{21}{1:07:01}: take the expectation of \eqref{ch15:eq-gbm} (the
$\dd W$ term has mean zero), $\dd\,\E[S_t]=\mu\,\E[S_t]\dd t$, so $\E[S_T]=S_te^{\mu(T-t)}$ for the real-world dynamics. In the risk-neutral
world the stock must grow at the risk-free rate, $\E_\Q[S_T]=S_te^{r(T-t)}$ (this was the property $qS_1+(1-q)S_2=S_0e^{r\dd t}$ of the tree), which forces $\mu\to r$.
The volatility is \emph{unchanged}: the size of the fluctuations is a pathwise property (the quadratic variation, \cref{ch11:thm-qv}), the same under every equivalent measure, whereas the drift
can be changed at will. Changing only the drift of a Brownian motion is the content of Girsanov's theorem (\cref{ch:sde}): $W^{\Q}_t=W_t+\lambda t$ with
$\lambda=(\mu-r)/\sigma$, the \emph{market price of risk} (the Sharpe ratio of the stock; this identification is not made in the lecture). In this language
\begin{equation}\label{ch15:eq-girsanov}
  \dd S=\mu S\dd t+\sigma S\dd W=rS\dd t+\sigma S\bigl(\dd W+\lambda\dd t\bigr)=rS\dd t+\sigma S\dd W^{\Q}.
\end{equation}
The lognormal density of the terminal price under $\Q$ is (slide~21)
\begin{equation}\label{ch15:eq-density}
  p_\Q(S_T\mid S_t)=\frac1{S_T\,\sigma\sqrt{2\pi(T-t)}}\exp\left[-\frac{\bigl(\ln(S_T/S_t)-(r-\sigma^2/2)(T-t)\bigr)^2}{2\sigma^2(T-t)}\right],
\end{equation}
the density of the lognormal variable \cref{ch11:eq-gbm} with $\mu\to r$ \ts{24}{21}{1:09:15}. Pricing a claim then means ``integrate the payoff against this density and discount'', instead of solving
a PDE \ts{24}{21}{1:09:15}; this is easy for a lognormal density and typically not for other models. The equivalence of \eqref{ch15:eq-rnval} with the PDE is the Feynman--Kac
formula (\cref{ch:sde}), already visible in \eqref{ch15:eq-pde-solution}.

\begin{coursenote}[Erratum: the density on slide 21 (2024) and 20 (2013)]
The pre-factor of the printed density is wrong in both years. In 2024 it reads $S_t/(\sigma S_T\sqrt{2\pi T})$, in 2013 $1/(\sigma S\sqrt{2\pi T})$. The correct normalisation, which makes the density integrate to one
in $S_T$, is $1/\bigl(S_T\,\sigma\sqrt{2\pi(T-t)}\bigr)$, as in \eqref{ch15:eq-density}: the $S_T$ in the denominator is the Jacobian of $S_T=e^{X}$ and the time is the \emph{time to expiry}
$T-t$, as in the exponent, not $T$. The exponent on the slide is correct.
\end{coursenote}

\begin{intuition}[Why ``neutral'']
There is no claim that investors are risk-neutral. Real-world prices already contain risk premia: $\mu>r$ is the reward for bearing the risk of the stock. The claim is that the \emph{price of a derivative relative to the
stock} does not depend on those premia because the derivative can be replicated by the stock and the bank; so one may compute the price in any world where
the stock earns $r$, and the simplest is the one where nobody cares about risk. Changing $\mu\to r$ moves the drift of the log-price, not its width: a physicist would say that $\Q$ is the same
diffusion in a different, tilted, ``potential'' (the tilt is the drift), and the Radon--Nikodym factor $\exp(-\lambda W_T-\tfrac12\lambda^2T)$ is a
Boltzmann-like reweighting of the paths.
\end{intuition}

\begin{example}[The wrong expectation]\label{ch15:ex-wrong}
For $S_0=K=80$, $T=2$, $r=5\%$, $\sigma=20\%$ the correct price is $12.90$. Suppose the stock has real-world drift $\mu=12\%$. Discounting the \emph{real-world} expectation
of the payoff at the risk-free rate gives $e^{-rT}\E[(S_T-K)^+]=22.17$, a $72\%$ overestimate: it is not an arbitrage price, because one cannot replicate the option with it. A Monte Carlo estimate under $\Q$ with $2\cdot10^5$
paths gives $12.955\pm0.040$ (one standard error), consistent with $12.901$. (Computed for these notes.)
\end{example}

\begin{finance}[Where the approach is used]
For a desk the two languages are complementary. The PDE (finite-difference grids) is the tool for low-dimensional problems with early exercise (American options), barriers and local-vol models; the expectation
(Monte Carlo) for high-dimensional or path-dependent payoffs; trees are equivalent to an explicit finite-difference scheme (\cref{ch15:prop-limit}(ii)) \ts{24}{21}{1:04:56}\ts{13}{19}{34:59}.
Slide~22 of 2024 lists what the whole business uses: PDE, numerical methods, stochastic calculus, simulation, statistics, ``much, much more''.
\end{finance}

% ==================================================================
\section{The Black--Scholes formula}\label{ch15:sec-formula}

\begin{theorem}[Black--Scholes formula]\label{ch15:thm-bs}
In the model \eqref{ch15:eq-gbm}, with $\tau=T-t$, the European call and put are worth
\begin{align}
  C(S,t)&=S\,\Phi(d_1)-Ke^{-r\tau}\,\Phi(d_2)=e^{-r\tau}\bigl[e^{r\tau}S\,\Phi(d_1)-K\,\Phi(d_2)\bigr],\label{ch15:eq-call}\\
  P(S,t)&=Ke^{-r\tau}\,\Phi(-d_2)-S\,\Phi(-d_1)=e^{-r\tau}\bigl[K\,\Phi(-d_2)-e^{r\tau}S\,\Phi(-d_1)\bigr],\label{ch15:eq-put}
\end{align}
where $\Phi$ is the standard normal cdf and
\begin{equation}\label{ch15:eq-d12}
  d_1=\frac{\ln(S/K)+(r+\tfrac12\sigma^2)\tau}{\sigma\sqrt\tau},\qquad d_2=d_1-\sigma\sqrt\tau=\frac{\ln(S/K)+(r-\tfrac12\sigma^2)\tau}{\sigma\sqrt\tau}.
\end{equation}
\end{theorem}
The right-hand forms are the ones on slide~20 \ts{24}{21}{1:04:56}\ts{13}{19}{32:52}. In the lecture the formula is quoted as ``the solution of the PDE in terms of the error function, which has to be tabulated''; here are two derivations.

\begin{proof}[Proof 1: risk-neutral expectation]
By \eqref{ch15:eq-rnval}, $C=e^{-r\tau}\E[(S_T-K)^+]$ with $S_T=Se^{(r-\sigma^2/2)\tau+\sigma\sqrt\tau Z}$, $Z\sim N(0,1)$ (the marginal of \eqref{ch15:eq-density}). The event $\{S_T>K\}$ is
$\{Z>-d_2\}$ because $\ln(S_T/K)=\sigma\sqrt\tau\,(Z+d_2)$. So
\[
  \E\bigl[(S_T-K)^+\bigr]=\E\bigl[S_T\ind\{Z>-d_2\}\bigr]-K\,\Prob(Z>-d_2),\qquad\Prob(Z>-d_2)=\Phi(d_2).
\]
For the first term use the identity, for real $b,c$,
\[
  \E\bigl[e^{bZ}\ind\{Z>c\}\bigr]=\int_c^\infty e^{bz}\frac{e^{-z^2/2}}{\sqrt{2\pi}}\dd z=e^{b^2/2}\int_c^\infty\frac{e^{-(z-b)^2/2}}{\sqrt{2\pi}}\dd z=e^{b^2/2}\,\Phi(b-c)
\]
(completing the square). With $b=\sigma\sqrt\tau$, $c=-d_2$ one has $b-c=d_2+\sigma\sqrt\tau=d_1$ and
$\E[S_T\ind\{Z>-d_2\}]=Se^{(r-\sigma^2/2)\tau}e^{\sigma^2\tau/2}\Phi(d_1)=Se^{r\tau}\Phi(d_1)$. Multiplying by $e^{-r\tau}$ gives \eqref{ch15:eq-call}. The put follows in the same way with
$\{Z<-d_2\}$, or from \cref{ch15:thm-parity} below.
\end{proof}
This is also \cref{ch03:cor-normal}(ii) with $\mu_{\ln}=\ln S+(r-\tfrac12\sigma^2)\tau$ and $\sigma_{\ln}=\sigma\sqrt\tau$, exactly as anticipated in the finance box of that chapter; and it is the
two-line computation suggested in 2013 as a final-paper topic \ts{13}{19}{34:59}.

\begin{proof}[Proof 2: solving the PDE]
By \eqref{ch15:eq-pde-solution} the solution of \eqref{ch15:eq-bs} with final condition $(S-K)^+$ and the boundary conditions \eqref{ch15:eq-bccall} is the same Gaussian integral as in Proof~1; uniqueness of
the bounded-growth solution of the heat equation (\cref{ch11:prop-heat}) shows that it is \emph{the} solution. One can also verify directly that \eqref{ch15:eq-call} solves \eqref{ch15:eq-bs}, using the identity
$S\varphi(d_1)=Ke^{-r\tau}\varphi(d_2)$ (\cref{ch15:sec-greeks}), and that $C\to(S-K)^+$ as $\tau\to0$ (for $S>K$, $d_{1,2}\to+\infty$; for $S<K$, $d_{1,2}\to-\infty$).
\end{proof}

\begin{example}[The at-the-money call of the slides]\label{ch15:ex-atm}
For $S=K=80$, $\tau=2$, $r=5\%$, $\sigma=20\%$: $d_1=0.14/0.2828=0.4950$, $d_2=0.2121$, $\Phi(d_1)=0.6897$, $\Phi(d_2)=0.5840$,
$S\Phi(d_1)=55.175$, $Ke^{-r\tau}\Phi(d_2)=72.39\cdot0.5840=42.274$ and $C=12.901$. By parity the put is $5.288$.
The two terms are the price of \emph{receiving the stock if $S_T>K$} and of \emph{paying $K$ if $S_T>K$}, the number $\Phi(d_2)=0.584$ being the risk-neutral probability of finishing in the money.
\end{example}

\begin{coursenote}[On the parameters of the slides]
The slides label the plots with $S=80$, $K=80$, $T=2$ years (forward: $K=88.41$) and never state $r$ or $\sigma$. The forward strike identifies $r=5\%$ (continuously compounded, $80e^{0.1}=88.41$), and the digital of slide~26 tends to $0.905=e^{-0.1}$ for large $S$ and takes the value $\approx0.53$ at $S=80$;
with $\sigma=20\%$ formula \eqref{ch15:eq-digital} gives $0.528$. All figures in this chapter that reproduce the slides use $r=5\%$, $\sigma=20\%$; this identification is ours.
\end{coursenote}

\begin{intuition}[Reading the formula]
\begin{itemize}[nosep]
  \item $e^{-r\tau}\Phi(d_2)=$ price of the \emph{cash-or-nothing} (digital) option, $\Phi(d_2)=\Q(S_T>K)$; $S\Phi(d_1)$ is the price of the \emph{asset-or-nothing} option, which pays $S_T$ if $S_T>K$.
        $\Phi(d_1)$ is a probability too, but under the measure in which the stock, not the bank, is the numeraire: the probability that the option is exercised as seen by a holder of the stock.
  \item As $\sigma\sqrt\tau\to0$ the option becomes the forward $(S-Ke^{-r\tau})^+$: the time value is the smoothing of the hockey stick over a width $\sigma\sqrt\tau$ (\cref{ch15:fig-payoffs}).
  \item At the money forward ($Se^{r\tau}=K$), $d_1=-d_2=\tfrac12\sigma\sqrt\tau$ and $C=S\,[2\Phi(\tfrac12\sigma\sqrt\tau)-1]\approx0.4\,S\sigma\sqrt\tau$, the Bachelier rule of thumb quoted in \cref{ch:probability} (not discussed in class).
  \item The price depends on $S,K,\tau,r,\sigma$ but not on $\mu$. The only quantity that is not observed is $\sigma$, which is what traders quote instead of a price (\cref{ch15:sec-smile}).
\end{itemize}
\end{intuition}

\begin{exercise}[not from the course: the formula by expectation]\label{ch15:ex-formula}
Derive the put formula \eqref{ch15:eq-put} directly from $P=e^{-r\tau}\E_\Q[(K-S_T)^+]$ and check parity. Evaluate the call and the put for $S=100$, $K=90$, $\tau=0.5$, $r=3\%$, $\sigma=25\%$.
Show that $C$ is increasing and convex in $S$, increasing in $\sigma$ and $\tau$, and that $(S-Ke^{-r\tau})^+\le C\le S$ (the no-arbitrage bounds).
\end{exercise}

\begin{exercise}[not from the course: risk-neutral drift]\label{ch15:ex-girsanov}
Apply \eqref{ch15:eq-ito} to $e^{-rt}S_t$ to show that $\dd(e^{-rt}S_t)=(\mu-r)e^{-rt}S_t\dd t+\sigma e^{-rt}S_t\dd W$ and, with $W^{\Q}_t=W_t+\lambda t$, that it is a martingale for $\lambda=(\mu-r)/\sigma$. Compute
$\E_{\Prob}[C_T]$ and $\E_{\Q}[C_T]$ for the call of \cref{ch15:ex-atm} with $\mu=12\%$ and explain why $e^{-rT}\E_{\Prob}[C_T]$ is not the price.
\end{exercise}

% ==================================================================
\section{Put--call parity: a model-free replication}\label{ch15:sec-parity}
The lecture illustrates replication with a relation that requires \emph{no} assumption on the stock dynamics \ts{24}{21}{1:10:22}\ts{13}{19}{38:07}.

\begin{theorem}[Put--call parity]\label{ch15:thm-parity}
For European options on a non-dividend-paying asset with the same strike $K$ and expiry $T$, in the absence of arbitrage,
\begin{equation}\label{ch15:eq-parity}
  C(t)-P(t)=S_t-Ke^{-r(T-t)}.
\end{equation}
\end{theorem}
\begin{proof}
The payoffs at $T$ satisfy $(S_T-K)^+-(K-S_T)^+=S_T-K$ (the ``hockey stick'' and its reflection add up to a straight line \ts{24}{21}{1:11:25}). So buying a call and selling a put
replicates a forward with strike $K$ \emph{statically}, i.e.\ with no trading before $T$: it is a forward contract, and by
\cref{ch15:prop-fwd} (applied over $[t,T]$) it is worth $S_t-Ke^{-r(T-t)}$ (stock, minus the bond that pays $K$). The two portfolios have the same payoff in every state, so they have the same price. \ts{24}{21}{1:12:31}
\end{proof}
If \eqref{ch15:eq-parity} failed, say $C-P>S-Ke^{-r\tau}$, one would sell the call, buy the put and the stock, and borrow $Ke^{-r\tau}$, locking in the difference at no risk. The remark of the lecture is that the relation is
\emph{independent of the dynamics}: one can price a put from the observed calls without a model, which matters in the real market where the stock is \emph{not} lognormal
\ts{24}{21}{1:12:31}\ts{13}{19}{38:07}. It is also the reason why the implied volatility of a call and of a put with the same strike and expiry must coincide.

\subsection{Parity on Bloomberg screens}
The slides check the relation on real quotes (mid $=$ average of bid and ask). The tables below use numbers read from the screens for these notes; a reading error of a cent is possible.
\begin{table}[ht]
\centering\small
\begin{tabular}{r r r r r r}
\toprule
$K$ & call mid & put mid & $C-P$ & $S-Ke^{-r\tau}$ & difference\\
\midrule
75 & 6.900 & 0.300 & 6.600 & 6.555 & $+0.045$\\
80 & 2.950 & 1.300 & 1.650 & 1.583 & $+0.067$\\
85 & 0.750 & 4.150 & $-3.400$ & $-3.390$ & $-0.010$\\
90 & 0.125 & 8.800 & $-8.675$ & $-8.362$ & $-0.313$\\
\bottomrule
\end{tabular}
\caption{IBM options expiring 22 April 2006, screen of 8 March 2006 (2013 slide~23), $S=81.14$; the discount uses $r=4.5\%$ and $\tau=45/365$ (our assumptions: the slide does not
show a rate, so this is a rough choice, and dividends are ignored). Equity options are American, and the deep in-the-money put ($K=90$, wide spread) carries an early-exercise premium that parity does not include.}\label{ch15:tab-ibm}
\end{table}

\begin{table}[ht]
\centering\small
\begin{tabular}{r r r r r r}
\toprule
$K$ & call mid & put mid & $C-P$ & $S-PV(D)-Ke^{-r\tau}$ & difference\\
\midrule
90 & 23.63 & 0.26 & 23.37 & 23.14 & $+0.23$\\
100 & 14.00 & 0.77 & 13.23 & 13.15 & $+0.08$\\
110 & 6.10 & 2.83 & 3.27 & 3.17 & $+0.10$\\
120 & 1.62 & 8.35 & $-6.73$ & $-6.81$ & $+0.08$\\
\bottomrule
\end{tabular}
\caption{AAPL options expiring 20 January 2017 ($\tau=81$ days), screen of 31 October 2016 (2024 slide~24), $S=113.54$, $r=0.83\%$ and the dividend of \$0.57 shown on the screen (its present value is
subtracted from $S$: parity for a dividend-paying stock is $C-P=S-PV(D)-Ke^{-r\tau}$, ``not discussed'' in class).}\label{ch15:tab-aapl}
\end{table}
On slide~24 the lecturer plots the call, the put and their difference against strike: the difference is a straight line of slope $-e^{-r\tau}\approx-1$, while the implied volatility
of the same options is a U-shaped curve (see \cref{ch15:sec-smile}). A straight line is what parity requires, and ``if it wouldn't be, people would jump on it immediately'' \ts{24}{21}{1:13:36}.
In 2013 the same idea is shown on the IBM screen of 2006 (slide~23), at the end of the lecture: far from the strike a call is worth $S-K$ ($S-K=81.14-55=26.14$, against a quoted mid of $26.50$), and between
strikes the prices change by 5 when the strike changes by 5 \ts{13}{19}{36:03}.

\begin{finance}[Replication as a way of thinking]
The lecturer's closing message \ts{24}{21}{1:14:40}\ts{13}{19}{48:38}: ``what you try to do is look at the payout of a portfolio and replicate it with the instruments at hand, or at least approximate it,
because if you did, you know the price now, or you know an arbitrage''. Conversely one builds a complicated portfolio of listed options and looks for a mispricing with respect to a replicating one: nowadays this
is ``practically impossible'' for vanilla options, but was done ``25 years ago'' with more complex options in FX (a story he leaves to J.~Xia). The hedge of a book of exotics is usually dynamic and needs rebalancing
\ts{13}{19}{49:42}; static replications like parity and the call spread below are the exceptions, and the most robust.
\end{finance}

\begin{exercise}[not from the course: put--call parity]\label{ch15:ex-parity}
Prove parity when the stock pays a known cash dividend $D$ at $t_D\in(t,T)$: $C-P=S-De^{-r(t_D-t)}-Ke^{-r\tau}$. Using the numbers of \cref{ch15:tab-ibm} choose an interest rate $r$ that best fits the four strikes by least squares, and comment on the residuals.
Design an arbitrage (list the trades and the guaranteed profit) if the call in the $K=100$ row of \cref{ch15:tab-aapl} were quoted at 14.50 instead of 14.00.
\end{exercise}

% ==================================================================
\section{Digital options and call spreads}\label{ch15:sec-digital}
The lecture ends with a challenge \ts{24}{21}{1:15:43}, which in 2013 was posed as an interview question \ts{13}{19}{43:24}: \emph{price and hedge a digital (all-or-nothing) option}, which pays $1$
if $S_T>K$ and $0$ otherwise, given that the prices of calls of \emph{all strikes} are known. Slide~26 of 2024 compares the Black--Scholes digital with a ``call spread'' built from the same
AAPL data.

\paragraph{Replication.} A call spread that is long a call with strike $K-\varepsilon$ and short a call with strike $K+\varepsilon$ pays
$(S_T-K+\varepsilon)^+-(S_T-K-\varepsilon)^+$, a ramp from $0$ (for $S_T\le K-\varepsilon$) to $2\varepsilon$ (for $S_T\ge K+\varepsilon$). A student proposes it and the lecturer confirms:
$\pm\tfrac12$ around $K$ gives a ramp up to $1$, $\pm\tfrac14$ needs two spreads, and in general the payoff of $\frac1{2\varepsilon}$ spreads is a ramp of unit height and width $2\varepsilon$
\ts{13}{19}{44:31}. Letting $\varepsilon\to0$, the digital price is
\begin{equation}\label{ch15:eq-digital-limit}
  D(K)=\lim_{\varepsilon\to0}\frac{C(K-\varepsilon)-C(K+\varepsilon)}{2\varepsilon}=-\frac{\partial C}{\partial K}
\end{equation}
\emph{whatever the model} \ts{13}{19}{46:37}, since $C(K)=e^{-r\tau}\E_\Q[(S_T-K)^+]$ and $\partial_K(S_T-K)^+=-\ind\{S_T>K\}$; this is the Breeden--Litzenberger relation of \cref{ch03:cor-bl} ($D=e^{-r\tau}\Q(S_T>K)$, while $\partial^2_KC=e^{-r\tau}p_\Q(K)$). The lecturer notes that the answer is a
derivative of the call price with respect to $K$, and that it should not be surprising: the derivative of a ramp is a step \ts{13}{19}{46:37}; the same idea works for any payoff, e.g.\ a box. In practice the dealer
sells the digital to a client and hedges with two liquid calls on the exchange, and digitals are less liquid than calls \ts{13}{19}{47:37}.

In Black--Scholes the digital costs
\begin{equation}\label{ch15:eq-digital}
  D=e^{-r\tau}\,\Phi(d_2),
\end{equation}
which for $S=K=80$, $\tau=2$, $r=5\%$, $\sigma=20\%$ is $0.5284$; the call spread with half-width $\varepsilon=10,5,2,1,0.1$ costs $0.5291,\ 0.5286,\ 0.52846,\ 0.52843,\ 0.528423$ (computed for these notes).
For large $S$, $D\to e^{-r\tau}=0.905$, the value of a certain \$1 at $T$, as in the left graph of slide~26 \ts{24}{21}{1:16:44}.

\begin{figure}[ht]
\centering
\begin{tikzpicture}
\begin{axis}[width=0.9\linewidth,height=5.6cm,axis lines=left,xmin=20,xmax=140,ymin=0,ymax=1.08,
  xlabel={$S$},tick label style={font=\footnotesize},label style={font=\small},
  legend style={at={(0.03,0.97)},anchor=north west,font=\scriptsize,draw=none,fill=none}]
  \addplot[thick,black!45] coordinates {(20,0) (80,0) (80,1) (140,1)};
  \addplot[very thick,defcol] coordinates {(20,0.0000) (23,0.0000) (26,0.0001) (29,0.0003) (32,0.0011) (35,0.0030) (38,0.0070) (41,0.0142) (44,0.0259) (47,0.0431) (50,0.0666) (53,0.0967) (56,0.1331) (59,0.1753) (62,0.2220) (65,0.2722) (68,0.3244) (71,0.3772) (74,0.4295) (77,0.4802) (80,0.5284) (83,0.5736) (86,0.6153) (89,0.6534) (92,0.6877) (95,0.7183) (98,0.7453) (101,0.7691) (104,0.7897) (107,0.8076) (110,0.8230) (113,0.8362) (116,0.8474) (119,0.8568) (122,0.8649) (125,0.8716) (128,0.8773) (131,0.8820) (134,0.8859) (137,0.8892) (140,0.8920)};
  \addplot[thick,dashed,thmcol] coordinates {(20,0.0000) (23,0.0000) (26,0.0001) (29,0.0005) (32,0.0016) (35,0.0040) (38,0.0088) (41,0.0172) (44,0.0302) (47,0.0489) (50,0.0738) (53,0.1050) (56,0.1421) (59,0.1844) (62,0.2309) (65,0.2803) (68,0.3313) (71,0.3827) (74,0.4334) (77,0.4824) (80,0.5291) (83,0.5729) (86,0.6134) (89,0.6504) (92,0.6838) (95,0.7138) (98,0.7405) (101,0.7640) (104,0.7846) (107,0.8025) (110,0.8180) (113,0.8314) (116,0.8429) (119,0.8527) (122,0.8610) (125,0.8681) (128,0.8741) (131,0.8791) (134,0.8834) (137,0.8869) (140,0.8899)};
  \legend{payoff $\ind\{S>K\}$,{Black--Scholes $D(t,S)$},call spread $\varepsilon=10$}
\end{axis}
\end{tikzpicture}
\caption{The digital option ($K=80$, $\tau=2$, $r=5\%$, $\sigma=20\%$) and the call spread with strikes $70$ and $90$. The spread is already close to the digital price; for $\varepsilon\to0$ it converges to it.}\label{ch15:fig-digital}
\end{figure}

\paragraph{With market data.} The right graph of slide~26 is the call spread built from the AAPL quotes (slide~24) as a function of $K-S$. From the numbers of \cref{ch15:tab-aapl}, $(C(110)-C(120))/10=0.448$
and $(C(105)-C(125))/20=0.451$ estimate the digital with $K=115$ (probability of finishing in the money, discounted). A flat-volatility Black--Scholes digital with the implied volatility of the $K=115$ call ($19.8\%$) would give only $0.413$: the
difference is caused by the \emph{skew} of the smile ($\partial\sigma/\partial K<0$ there) and is a first-principles example of ``the world is not Black--Scholes'' (\cref{ch15:sec-smile}). The calculation is ours (not in class).

\begin{intuition}[Static versus dynamic replication]
Any smooth payoff $g(S_T)$ is a continuous portfolio of calls, $g(S_T)=g(0)+g'(0)S_T+\int_0^\infty g''(K)(S_T-K)^+\dd K$ (a Taylor formula with integral remainder in $K$; not discussed in class), so its price
is a fixed combination of call prices with no model at all. The Black--Scholes dynamic hedge is what one needs when the instruments are \emph{not} available at all strikes.
\end{intuition}

\begin{exercise}[not from the course: digital options]\label{ch15:ex-digital}
\leavevmode
\begin{enumerate}[(a)]
  \item Prove $D=-\partial C/\partial K$ for any model where $C(K)=e^{-r\tau}\E_\Q[(S_T-K)^+]$.
  \item Show that a call spread with strikes $K$ and $K+\varepsilon$, scaled by $1/\varepsilon$, \emph{super-replicates} the digital,
      so that the dealer's ask for the digital is $\le[C(K)-C(K+\varepsilon)]/\varepsilon$, and find the sub-replicating spread giving a bid.
  \item In a model with a smile, $\sigma=\sigma(K)$, show that $D=e^{-r\tau}\Phi(d_2)-\mathcal V\,\sigma'(K)$ and explain the sign
      of the correction for the AAPL data of \cref{ch15:tab-smile}.
\end{enumerate}
\end{exercise}


% ==================================================================
\section{Greeks and hedging}\label{ch15:sec-greeks}
The replication argument produced one hedge ratio, $a=\partial f/\partial S$ \eqref{ch15:eq-hedge}. Traders call the derivatives of the price with respect to its inputs the \emph{Greeks}; they measure
the risks of a position (``the mathematical apparatus allows one to compute the current price of a derivative and its risks'', slide~7). The lecture only uses the delta; the rest of this section is not
discussed in class but follows from \cref{ch15:thm-bs} and completes the picture.

\begin{proposition}[Greeks of the European call]\label{ch15:prop-greeks}
With $\varphi=\Phi'$ the standard normal density and $\tau=T-t$,
\begin{equation}\label{ch15:eq-greeks}
\begin{gathered}
  \Delta=\frac{\partial C}{\partial S}=\Phi(d_1),\qquad
  \Gamma=\frac{\partial^2C}{\partial S^2}=\frac{\varphi(d_1)}{S\sigma\sqrt\tau},\qquad
  \mathcal V=\frac{\partial C}{\partial\sigma}=S\varphi(d_1)\sqrt\tau,\\
  \Theta=\frac{\partial C}{\partial t}=-\frac{S\varphi(d_1)\sigma}{2\sqrt\tau}-rKe^{-r\tau}\Phi(d_2),\qquad
  \rho=\frac{\partial C}{\partial r}=K\tau e^{-r\tau}\Phi(d_2).
\end{gathered}
\end{equation}
For the put, $\Delta_P=\Delta-1$ and $\Gamma,\mathcal V$ are the same; $\Theta_P=-\frac{S\varphi(d_1)\sigma}{2\sqrt\tau}+rKe^{-r\tau}\Phi(-d_2)$ and $\rho_P=-K\tau e^{-r\tau}\Phi(-d_2)$.
The price satisfies \eqref{ch15:eq-bs} in the form $\Theta+\tfrac12\sigma^2S^2\Gamma+rS\Delta=rC$.
\end{proposition}
\begin{proof}
First the identity $S\varphi(d_1)=Ke^{-r\tau}\varphi(d_2)$: since $d_2=d_1-\sigma\sqrt\tau$,
$\varphi(d_2)=\varphi(d_1)\exp(\sigma\sqrt\tau\,d_1-\tfrac12\sigma^2\tau)$ and $\sigma\sqrt\tau\,d_1=\ln(S/K)+(r+\tfrac12\sigma^2)\tau$, so $\varphi(d_2)=\varphi(d_1)\,(S/K)\,e^{r\tau}$.
Now $\partial_Sd_1=\partial_Sd_2=1/(S\sigma\sqrt\tau)$ and
$\partial_SC=\Phi(d_1)+\bigl[S\varphi(d_1)-Ke^{-r\tau}\varphi(d_2)\bigr]\partial_Sd_1=\Phi(d_1)$, the bracket vanishing by the identity. Differentiating once more gives $\Gamma$.
For $\mathcal V$: $\partial_\sigma C=S\varphi(d_1)\partial_\sigma d_1-Ke^{-r\tau}\varphi(d_2)\partial_\sigma d_2=S\varphi(d_1)\,\partial_\sigma(d_1-d_2)=S\varphi(d_1)\sqrt\tau$, by the identity and $d_1-d_2=\sigma\sqrt\tau$. The theta and rho are computed in the same way (or theta from the PDE, once $\Delta,\Gamma$ are known);
the put values follow from parity $P=C-S+Ke^{-r\tau}$.
\end{proof}

Numbers for $S=K=100$, $\tau=1$, $r=5\%$, $\sigma=20\%$ ($d_1=0.35$, $d_2=0.15$; call $10.4506$, put $5.5735$; checked against finite differences and against the PDE, residual $10^{-16}$, for these notes):
\begin{center}\small
\begin{tabular}{l r r r r r}
\toprule
 & $\Delta$ & $\Gamma$ & $\mathcal V$ (per unit $\sigma$) & $\Theta$ (per year) & $\rho$ (per unit $r$)\\
\midrule
call & $0.6368$ & $0.01876$ & $37.52$ & $-6.414$ & $53.23$\\
put  & $-0.3632$ & $0.01876$ & $37.52$ & $-1.658$ & $-41.89$\\
\bottomrule
\end{tabular}
\end{center}
A one-point (1\%) rise in volatility adds about $0.375$ to the price, and one day of time decay of the call costs about $6.41/365=0.018$.

\begin{figure}[ht]
\centering
\begin{tikzpicture}
\begin{axis}[name=A,width=0.48\linewidth,height=5.2cm,xmin=60,xmax=140,ymin=0,ymax=1.05,axis lines=left,xlabel={$S$},
  tick label style={font=\footnotesize},label style={font=\small},title={\small delta $\Phi(d_1)$},title style={yshift=-3pt},
  legend style={at={(0.03,0.97)},anchor=north west,font=\scriptsize,draw=none,fill=none}]
  \addplot[very thick,defcol] coordinates {(60,0.0138) (62,0.0207) (64,0.0300) (66,0.0420) (68,0.0572) (70,0.0759) (72,0.0981) (74,0.1239) (76,0.1533) (78,0.1861) (80,0.2219) (82,0.2604) (84,0.3009) (86,0.3431) (88,0.3862) (90,0.4298) (92,0.4733) (94,0.5162) (96,0.5580) (98,0.5983) (100,0.6368) (102,0.6733) (104,0.7075) (106,0.7394) (108,0.7688) (110,0.7958) (112,0.8203) (114,0.8426) (116,0.8626) (118,0.8805) (120,0.8965) (122,0.9106) (124,0.9230) (126,0.9339) (128,0.9434) (130,0.9517) (132,0.9589) (134,0.9651) (136,0.9704) (138,0.9750) (140,0.9789)};
  \addplot[very thick,thmcol,dashed] coordinates {(60,0.0000) (62,0.0000) (64,0.0000) (66,0.0000) (68,0.0001) (70,0.0003) (72,0.0009) (74,0.0023) (76,0.0051) (78,0.0105) (80,0.0199) (82,0.0352) (84,0.0584) (86,0.0912) (88,0.1349) (90,0.1898) (92,0.2550) (94,0.3286) (96,0.4078) (98,0.4892) (100,0.5695) (102,0.6454) (104,0.7147) (106,0.7757) (108,0.8276) (110,0.8704) (112,0.9046) (114,0.9313) (116,0.9515) (118,0.9664) (120,0.9772) (122,0.9847) (124,0.9900) (126,0.9935) (128,0.9959) (130,0.9974) (132,0.9984) (134,0.9990) (136,0.9994) (138,0.9997) (140,0.9998)};
  \legend{$\tau=1$,$\tau=0.25$}
\end{axis}
\begin{axis}[at={(A.east)},anchor=west,xshift=1.1cm,width=0.48\linewidth,height=5.2cm,xmin=60,xmax=140,ymin=0,ymax=0.07,axis lines=left,xlabel={$S$},
  tick label style={font=\footnotesize},label style={font=\small},title={\small gamma},title style={yshift=-3pt},
  legend style={at={(0.03,0.97)},anchor=north west,font=\scriptsize,draw=none,fill=none}]
  \addplot[very thick,defcol] coordinates {(60,0.0029) (62,0.0040) (64,0.0053) (66,0.0068) (68,0.0084) (70,0.0102) (72,0.0120) (74,0.0138) (76,0.0156) (78,0.0172) (80,0.0186) (82,0.0198) (84,0.0207) (86,0.0214) (88,0.0217) (90,0.0218) (92,0.0216) (94,0.0212) (96,0.0206) (98,0.0197) (100,0.0188) (102,0.0177) (104,0.0165) (106,0.0153) (108,0.0141) (110,0.0129) (112,0.0117) (114,0.0106) (116,0.0095) (118,0.0085) (120,0.0075) (122,0.0066) (124,0.0058) (126,0.0051) (128,0.0044) (130,0.0039) (132,0.0033) (134,0.0029) (136,0.0025) (138,0.0021) (140,0.0018)};
  \addplot[very thick,thmcol,dashed] coordinates {(60,0.0000) (62,0.0000) (64,0.0000) (66,0.0000) (68,0.0001) (70,0.0002) (72,0.0004) (74,0.0010) (76,0.0019) (78,0.0036) (80,0.0060) (82,0.0095) (84,0.0139) (86,0.0191) (88,0.0247) (90,0.0301) (92,0.0349) (94,0.0385) (96,0.0404) (98,0.0407) (100,0.0393) (102,0.0365) (104,0.0327) (106,0.0282) (108,0.0236) (110,0.0192) (112,0.0151) (114,0.0116) (116,0.0087) (118,0.0063) (120,0.0045) (122,0.0031) (124,0.0022) (126,0.0014) (128,0.0009) (130,0.0006) (132,0.0004) (134,0.0002) (136,0.0001) (138,0.0001) (140,0.0001)};
  \legend{$\tau=1$,$\tau=0.25$}
\end{axis}
\end{tikzpicture}
\caption{Delta and gamma of a call with $K=100$, $r=5\%$, $\sigma=20\%$. Near expiry delta approaches the step function of the payoff and gamma concentrates around the strike, which makes short-dated at-the-money options the
hardest to hedge.}\label{ch15:fig-greeks}
\end{figure}

\paragraph{Delta hedging.} To be short one call and hedged one holds $\Delta=\Phi(d_1)$ shares (about $64$ shares per $100$ calls in the example above) and a negative cash position; the combination has
no first-order sensitivity to $S$. The residual risk is second order: over a small move $\delta S$ the hedged book changes by $\Theta\,\delta t+\tfrac12\Gamma\,\delta S^2$,
and the PDE $\Theta+\tfrac12\sigma^2S^2\Gamma=r(C-S\Delta)$ says that, on average, the gain from gamma ($\tfrac12\Gamma\,\delta S^2\approx\tfrac12\Gamma\sigma^2S^2\delta t$) pays for the time decay: a long-gamma
position bleeds theta and earns it back through movement, exactly at the rate that makes the hedged book earn $r$. This is the origin of the trader's statement that an option trade is a bet on realised versus implied volatility (\cref{ch:volatility}).

\begin{table}[ht]
\centering\small
\begin{tabular}{r r r r}
\toprule
rebalancings $n$ & mean error & std.\ dev. & std.\ dev.\ $\times\sqrt n$\\
\midrule
1 & $-0.208$ & 6.109 & 6.11\\
2 & $-0.088$ & 4.379 & 6.19\\
4 & $-0.092$ & 3.225 & 6.45\\
12 & $-0.040$ & 1.902 & 6.59\\
52 & $-0.014$ & 0.927 & 6.69\\
252 & $-0.005$ & 0.426 & 6.76\\
\bottomrule
\end{tabular}
\caption{Error of discrete delta hedging of a short call ($S_0=K=100$, $T=1$, $r=5\%$, $\sigma=20\%$, price $10.4506$ received) in a market whose real-world drift is $\mu=10\%\ne r$:
terminal value of the hedge portfolio minus the payoff, $2\cdot10^4$ simulated paths, hedge ratio $\Phi(d_1)$ recomputed at $n$ equally spaced dates. Computed for these notes.}\label{ch15:tab-hedge}
\end{table}

\paragraph{What discrete rebalancing costs.} The derivation \eqref{ch15:eq-selffin} assumed that the hedge is adjusted \emph{continuously}. \Cref{ch15:tab-hedge} shows what happens with $n$ rebalancings: the hedge error
has a small mean (below $0.21$ here, and shrinking with $n$) and standard deviation proportional to $n^{-1/2}$ (the last column is roughly constant), i.e.\ proportional to $\sqrt{\dd t}$: the risk is
not eliminated for a finite $n$, only made small, and every rebalancing has a transaction cost, so the desk chooses a finite frequency \ts{13}{19}{28:31}\ (``every time $\dd t$ we will have to rebalance'').

\begin{finance}[What breaks the assumptions]
Discrete hedging and transaction costs (above); jumps and gaps, for which no hedge in the stock alone is perfect and the market is incomplete (\cref{sec:oneperiod}: a non-replicable claim has an interval of arbitrage-free prices)
\ts{24}{21}{1:14:40}; stochastic volatility, so that $\sigma$ itself needs hedging with other options; funding and credit spreads (\cref{ch:counterparty}); dividends and borrow costs. Each is a chapter of its own in a derivatives library;
Black--Scholes remains the language in which prices are \emph{quoted}.
\end{finance}

\begin{exercise}[not from the course: Greeks]\label{ch15:ex-greeks}
Prove $S\varphi(d_1)=Ke^{-r\tau}\varphi(d_2)$, derive $\Delta,\Gamma,\mathcal V,\Theta$ and check that they satisfy the Black--Scholes PDE. For an at-the-money-forward option show that $C\approx0.4\,S\sigma\sqrt\tau$; prove in general that $\mathcal V=S^2\sigma\tau\,\Gamma$. What is the $\Delta$ of a call with $S$ very large? Of a call very close to expiry, at the money?
\end{exercise}

\begin{exercise}[not from the course: delta hedging by simulation]\label{ch15:ex-hedge}
Reproduce \cref{ch15:tab-hedge}: simulate GBM with drift $\mu$, sell a call at the Black--Scholes price, rebalance the delta at $n$ dates, and record the terminal error. Check that the distribution does not depend on $\mu$ (for large $n$),
that its width scales as $n^{-1/2}$, and add a proportional transaction cost $c|\Delta a|S$ to find the $n$ that minimises the cost-plus-risk for your favourite risk penalty. What changes if you hedge with a wrong $\sigma$ (say $15\%$ when the true one is $20\%$)?
\end{exercise}

% ==================================================================
\section{Implied volatility and the smile}\label{ch15:sec-smile}
The formula has one input that is not observed, $\sigma$. Since $\partial C/\partial\sigma=\mathcal V>0$ and $C\to(S-Ke^{-r\tau})^+$ as $\sigma\to0$, $C\to S$ as $\sigma\to\infty$, the map $\sigma\mapsto C_{BS}(\sigma)$ is a strictly increasing
bijection from $(0,\infty)$ onto $\bigl((S-Ke^{-r\tau})^+,S\bigr)$; hence every market price in that interval determines a unique \emph{implied volatility} $\sigma_{\mathrm{imp}}$ (found numerically, e.g.\ by Newton's method with the vega as derivative). Parity implies that call and put
of the same strike and expiry have the same $\sigma_{\mathrm{imp}}$.

In Black--Scholes $\sigma$ is one constant for the stock, so the implied volatility should be the \emph{same for all strikes and expiries}. The lecture points out that it is not
\ts{24}{21}{1:13:36}\ts{13}{19}{37:04}: plotted against strike it is a curve, a ``smile'' (in 2013 the word is ``skewed, or rather smiled''), shown in the lower right graph of slide~24. \Cref{ch15:tab-smile} and
\cref{ch15:fig-smile} report the numbers of that screen.

\begin{table}[ht]
\centering\small
\begin{tabular}{r r r r}
\toprule
$K$ & call mid & screen implied vol.\ (\%) & recomputed (\%)\\
\midrule
90 & 23.63 & 32.53 & 33.12\\
100 & 14.00 & 23.71 & 24.13\\
110 & 6.10 & 20.47 & 20.69\\
120 & 1.62 & 19.06 & 19.08\\
130 & 0.31 & 19.41 & 19.37\\
140 & 0.09 & 21.45 & 21.79\\
\bottomrule
\end{tabular}
\caption{AAPL calls of \cref{ch15:tab-aapl} ($S=113.54$, $\tau=81$ days). ``Recomputed'' inverts the European Black--Scholes formula with $S-PV(D)$, $r=0.83\%$ (our computation); the small differences from the Bloomberg values reflect the
American-style and dividend conventions of the screen, and are largest for deep in-the-money and far out-of-the-money quotes.}\label{ch15:tab-smile}
\end{table}

\begin{figure}[ht]
\centering
\begin{tikzpicture}
\begin{axis}[width=0.85\linewidth,height=5.6cm,axis lines=left,xmin=70,xmax=150,ymin=15,ymax=55,
  xlabel={strike $K$ ($S=113.54$)},ylabel={implied volatility (\%)},tick label style={font=\small},label style={font=\small}]
  \addplot[very thick,defcol,mark=*,mark size=1.4pt] coordinates {(75,51.08) (80,45.65) (82.5,43.27) (85,38.82) (90,32.53) (92.5,29.02) (95,26.92) (97.5,25.21) (100,23.71) (105,22.04) (110,20.47) (115,19.57) (120,19.06) (125,18.91) (130,19.41) (135,20.25) (140,21.45) (145,22.37)};
  \draw[thmcol,dashed] (axis cs:113.54,15) -- (axis cs:113.54,55) node[pos=0.93,right,font=\scriptsize]{spot};
\end{axis}
\end{tikzpicture}
\caption{Implied volatility of the AAPL calls, screen of 31 October 2016 (values read from 2024 slide~24; the quote at $K=87.5$ is missing). A constant-volatility model would give a horizontal line.}\label{ch15:fig-smile}
\end{figure}

The curve decreases steeply for low strikes and turns up slightly for high ones: options in the wings are relatively \emph{expensive}. Two consequences. First, the lognormal distribution is not the distribution the market uses: returns have
heavier tails and skew (\cref{ch03:sec-heavy}), volatility moves with the price, and the wings price in crash risk. Second, prices of \emph{non-vanilla} claims cannot be read off one volatility number: as in the digital
of \cref{ch15:sec-digital}, the skew enters. The lecturer's conclusion \ts{24}{21}{1:14:40}\ is that Black--Scholes is a ``nice example'' of the method, but more complicated dynamics are needed for stocks: \emph{stochastic volatility}
(the volatility itself follows an SDE; \cref{ch:volatility,ch:sde}) or \emph{jump-diffusion}. With jumps ``life becomes more complicated'': the market is incomplete and martingale
techniques for continuous paths no longer suffice. The risk-neutral \emph{method} is unchanged: find the measure(s) $\Q$ under which discounted traded assets are martingales, and price by discounted expectation; what changes is that (for stochastic volatility and jumps) $\Q$
is no longer unique (\cref{thm:ftap}, part~2), and calibration to the observed option prices (the smile) is what selects it.

\begin{exercise}[not from the course: implied volatility]\label{ch15:ex-iv}
Implement Newton's method for $\sigma_{\mathrm{imp}}$ with the vega as derivative, starting from $\sigma_0=\sqrt{2|\ln(S/K)+r\tau|/\tau}$. Recompute the last column of \cref{ch15:tab-smile}. Why does the iteration fail for deep out-of-the-money options and what is a safer alternative (bisection, Brent)?
\end{exercise}

% ==================================================================
\section{The two lectures compared, and errata}\label{ch15:sec-compare}
\begin{coursenote}[2013 versus 2024]
The core slides are the same: the 2024 deck is the 2013 deck (V.~Strela was the speaker in both) with five changes.
\begin{itemize}[nosep]
  \item \emph{Opening.} 2024 starts with the update of the yield curve (\cref{ch15:sec-curve}); 2013 goes straight to the bookie.
  \item \emph{More derivation.} The 2024 lecture spends about 15 minutes on why $\operatorname{sd}(\dd W)=\sqrt{\dd t}$ and on the Taylor expansion behind It\^o's formula and on a student's objection to $(\dd W)^2=\dd t$
        (\cref{ch15:sec-lognormal}); 2013 assumes the students know It\^o from the lecture that precedes it (``did you talk about It\^o already? OK, cool'').
  \item \emph{Discrete market.} 2013 keeps the number $b$ of bond units and $B_0$ separate (its slide~10 has $b=\frac{S_1f_2-S_2f_1}{(S_1-S_2)B_0e^{r\dd t}}$), 2024 uses the cash amount $B$ directly; the risk-neutral probability is called $q$ in 2013 and $p_{rn}$ in 2024.
  \item \emph{Examples.} 2024 replaces the March 2006 IBM screen (2013 slide~23) with two recent Bloomberg screens (AAPL 2016 with the implied-volatility smile, and IBM 2015), adds the put--call parity plot and the real-data
        call spread (slide~26); 2013 discusses parity and the smile only orally, at the end, and treats the digital as an interview question with the audience answering.
  \item \emph{Suggestions.} 2013 explicitly proposes as final-paper topics the reduction to the heat equation and the derivation of the formula as a lognormal expectation
        \ts{13}{19}{30:40}, both carried out above; 2024 announces the next two (guest) lectures, one of them by J.~Hull on machine learning (\cref{ch:ml}).
\end{itemize}
In 2013 the lecturer also remarks that the two-year rate was then about $32.5$ basis points, so that the $5\%$ of the example is illustrative and not realistic \ts{13}{19}{17:43}.
\end{coursenote}

\begin{coursenote}[Errata found in the material]
\begin{itemize}
  \item \emph{Density of $S_T$} (2024 slide~21, 2013 slide~20): the pre-factor should be $1/(S_T\sigma\sqrt{2\pi(T-t)})$; see \cref{ch15:sec-rn}.
  \item \emph{Boundary condition of the call} (slide~19, 2013 slide~18): $C(S,t)\to S-Ke^{-r(T-t)}$ as $S\to\infty$, not only $\cong S$; see \cref{ch15:sec-pde}.
  \item \emph{Spoken slips} (2013 transcript): the bookie's payout ``\$12.25'' is \$12\,500 (\cref{ch15:sec-bookie}). (2024 transcript): ``our strike has to be 0'' after the argument for zero interest rates means $K=S_0$ ($K=S_0e^{r\dd t}$ in general) \ts{24}{21}{22:20};
        ``$S_2$ less than 0'' in the call diagram is $S_2<K$ \ts{24}{21}{26:32}.
\end{itemize}
\end{coursenote}

% ==================================================================
\begin{keyformulas}
\textbf{One step and tree} (continuous compounding, up/down factors $u,d$):
\begin{gather*}
  q=\frac{e^{r\dd t}-d}{u-d},\quad qu+(1-q)d=e^{r\dd t},\quad f_0=e^{-r\dd t}\bigl[qf_u+(1-q)f_d\bigr],\quad a=\frac{f_u-f_d}{S_u-S_d},\\
  \text{CRR: }u=e^{\sigma\sqrt{\dd t}}=1/d,\quad q=\tfrac12+\tfrac{r-\sigma^2/2}{2\sigma}\sqrt{\dd t}+O(\dd t).
\end{gather*}
\textbf{Model and PDE:}
\begin{gather*}
  \dd S=\mu S\dd t+\sigma S\dd W,\qquad \dd f=\bigl(f_t+\mu Sf_S+\tfrac12\sigma^2S^2f_{SS}\bigr)\dd t+\sigma Sf_S\dd W,\\
  f_t+\tfrac12\sigma^2S^2f_{SS}+rSf_S-rf=0,\qquad f(S,T)=\text{payoff},\qquad a=f_S .
\end{gather*}
\textbf{Risk-neutral valuation:}
\[
  f_t=e^{-r(T-t)}\E_\Q[f_T],\qquad \dd S=rS\dd t+\sigma S\dd W^{\Q},\qquad W^{\Q}_t=W_t+\tfrac{\mu-r}{\sigma}t .
\]
\textbf{Black--Scholes formula} ($\tau=T-t$):
\begin{gather*}
  C=S\Phi(d_1)-Ke^{-r\tau}\Phi(d_2),\qquad P=Ke^{-r\tau}\Phi(-d_2)-S\Phi(-d_1),\\
  d_{1,2}=\frac{\ln(S/K)+(r\pm\frac12\sigma^2)\tau}{\sigma\sqrt\tau},\qquad C-P=S-Ke^{-r\tau}.
\end{gather*}
\textbf{Greeks:}
\begin{gather*}
  \Delta_C=\Phi(d_1),\quad\Gamma=\frac{\varphi(d_1)}{S\sigma\sqrt\tau},\quad\mathcal V=S\varphi(d_1)\sqrt\tau,\quad S\varphi(d_1)=Ke^{-r\tau}\varphi(d_2),\\
  \Theta+\tfrac12\sigma^2S^2\Gamma+rS\Delta=rC .
\end{gather*}
\textbf{Digital and call spread:} $D=e^{-r\tau}\Phi(d_2)=-\partial C/\partial K=\lim_{\varepsilon\to0}\dfrac{C(K-\varepsilon)-C(K+\varepsilon)}{2\varepsilon}$.
\end{keyformulas}
